\documentclass[UTF8,fontset=fandol,zihao=-4,openany,oneside]{ctexbook}
\usepackage[a4paper,margin=25mm,headheight=15pt]{geometry}
\usepackage{amsmath,amssymb,mathtools,booktabs,longtable,array,enumitem,needspace}
\usepackage{microtype}
\usepackage[unicode,colorlinks=true,linkcolor=blue!45!black,urlcolor=blue!55!black]{hyperref}
\usepackage{xcolor,bookmark,fancyhdr}
\allowdisplaybreaks[2]
\setlength{\emergencystretch}{3em}
\setlist{itemsep=.35em,topsep=.4em}
\renewcommand{\arraystretch}{1.3}
\pagestyle{fancy}\fancyhf{}
\fancyhead[L]{\small RP01：完整中文解析}
\fancyhead[R]{\small v0.1.0 · 2026-10-17}
\fancyfoot[C]{\thepage}
\hypersetup{pdftitle={RP01完整中文论文解析},pdfauthor={Qi Liu},pdfsubject={Companion to preprint v0.1.0}}
\title{RP01 完整中文论文解析\\\large 从群计数到偶数半维数等号群分类}
\author{Qi Liu\\\small AI辅助整理的中文阅读材料}
\date{对应预印本 v0.1.0\\2026年10月17日}
\begin{document}
\frontmatter\maketitle
\chapter*{版本与使用说明}\addcontentsline{toc}{chapter}{版本与使用说明}
本解析对应英文论文 \emph{Extremal Klein Four Counts in Semi-Extraspecial 2-Groups}，
预印本 v0.1.0（2026-10-17），文件为 \texttt{RP01\_preprint.pdf}。
对应英文稿为19页；定义、命题与公式编号沿用该版本。
解析的章号仅用于教学导航；原稿定位以文中“原稿”编号为准。
这里的解释、背景与展开计算不构成新的研究定理；有标识的教学补充亦不作为原创贡献。
数学探索、证明发展、计算检查、文献核对与文本生成有生成式AI参与，
包括OpenAI Codex/ChatGPT及记录中的Muse Spark 1.3。
材料交付不表示作者已经阅读或人工核验全部论证，也不是外部同行评审。
阅读本解析不需要答题、打卡或等待另一次对话。
引用键与英文稿一致，书目信息集中列在最后一章。
原创解析文字采用\href{https://creativecommons.org/licenses/by/4.0/}{CC BY 4.0}；构建与阅读代码MIT，MathJax另保留Apache-2.0。
英文论文：\href{https://doi.org/10.5281/zenodo.23450418}{doi:10.5281/zenodo.23450418}；配套资料：\href{https://doi.org/10.5281/zenodo.23450447}{doi:10.5281/zenodo.23450447}。

\chapter*{阅读路线与术语索引}\addcontentsline{toc}{chapter}{阅读路线与术语索引}


本书不是论文摘要。第1章补足本题所需背景，第2—8章沿英文正文展开证明，最后解释两个保留的附录。可以从头读，也可以按英文编号跳转。无需回答问题或等待对话；不理解某一步时，将英文编号、中文节号或HTML链接带到另一个同项目对话，并让它先读\texttt{READING\_START.md}。



整条证明只有一条主线：\textbf{用平方与交换关系建立计数接口 → 用三个字符和及其对偶求界 → 从下端等号确定极化结构 → 保留全部线性细化 → 以同一个目标换基归一 → 转回群同构。} 数值界的端点构造独立完成，不调用后面的分类定理来证明自身可达。



\begin{longtable}{@{}p{0.3033\linewidth}@{\hspace{.025\linewidth}}p{0.3033\linewidth}@{\hspace{.025\linewidth}}p{0.3033\linewidth}@{}}
\toprule
要读的原稿位置 & 解析入口 & 这一步解决的问题 \\
\midrule
\endhead
定理1.1、1.2；§1 & \hyperlink{chapter.1}{第1章} & 准确对象、数值与分类范围 \\
命题2.1、引理2.2 & \hyperlink{section.2.1}{2.1—2.4} & 群与完整平方映射、严格等价、中心积 \\
引理2.3 & \hyperlink{section.2.5}{2.5—2.7} & 每个有序对和子群的倍率 \\
命题3.2 & \hyperlink{section.3.3}{3.3—3.5} & 三阶Walsh矩与三对偶和 \\
引理3.3、命题3.4—3.5 & \hyperlink{section.3.6}{3.6—3.10} & 成对初等因子、亏损与充要等号条件 \\
§4 & \hyperlink{chapter.4}{第4章} & 每个允许参数的端点为何可达 \\
引理5.1—5.3 & \hyperlink{section.5.1}{5.1—5.5} & 局部环、全部细化与保留常数的相位 \\
命题5.4—5.5、注5.6 & \hyperlink{section.5.6}{5.6—5.8} & 96是坐标映射数；显式传递性与群模型 \\
引理6.1—6.3、注6.4 & \hyperlink{section.6.1}{6.1—6.7} & 核心与整个余块的符号、双曲归一及共同目标 \\
定理1.2完整证明 & \hyperlink{section.6.8}{6.8} & 必要性、充分性和整个群类型唯一 \\
例7.1、引理7.2、例7.3 & \hyperlink{chapter.7}{第7章} & 最大性与分块正号、偶数结论的范围 \\
§8、附录A/B & \hyperlink{chapter.8}{第8章} & 既有工具、实际文献代入及版本专属旁例 \\
\bottomrule
\end{longtable}


\Needspace{5\baselineskip}
术语第一次使用时仍会解释。下面的索引用于找位置，不要求预先记住。



\begin{longtable}{@{}p{0.3033\linewidth}@{\hspace{.025\linewidth}}p{0.3033\linewidth}@{\hspace{.025\linewidth}}p{0.3033\linewidth}@{}}
\toprule
记号或术语 & 在本题中表示什么 & 阅读位置 \\
\midrule
\endhead
special / extraspecial / semi-extraspecial & 三种不同强度的群结构条件 & \hyperlink{section.1.2}{1.2、1.4} \\
\(W,V,n\) & 中心；中心商；商的半维数，\(\dim V=2n\) & \hyperlink{section.1.2}{1.2—1.4} \\
\(Q,B\) & 平方映射及其极化，分别记录平方与交换 & \hyperlink{section.1.3}{1.3}、\hyperlink{section.2.1}{2.1} \\
\(I,s,v_4\) & 自逆集合、自逆数（含单位元）、实际四元子群数 & \hyperlink{section.1.1}{1.1} \\
\(S,z,T_0,J_2,L\) & 零点集、零点数、零和有序对、交换自逆有序对、全零平面数 & \hyperlink{section.2.5}{2.5—2.7} \\
\(q_1,q_2,q_3\) & 三个非零标量分量，\(q_3=q_1+q_2\) & \hyperlink{section.1.4}{1.4} \\
Walsh变换、Gauss和、Arf符号 & 加线性频率后的正负求和；零频求和；非退化二次型的正负型 & \hyperlink{section.3.1}{3.1—3.5} \\
\(h,N,M_n,\mathcal P\) & \(2^n,2^{2n}\)、基准计数、三个Walsh变换的第三混合矩 & \hyperlink{section.1.5}{1.5}、\hyperlink{section.3.4}{3.4} \\
\(D,C,\rho\) & 三实际对偶的函数和、其极化、极化的根维数 & \hyperlink{section.3.5}{3.5—3.8} \\
\(A,p(t)\) & 相对极化算子；\(t^2+t+1\) & \hyperlink{section.3.6}{3.6—3.8} \\
初等因子、循环模、socle & 多项式作用的基本块；一个向量生成的模；其中最小非零子模 & \hyperlink{section.3.7}{3.7} \\
\(K,R,\varepsilon\) & 四元域；长度2局部环；非零平方为零元 & \hyperlink{section.5.2}{5.2—5.3} \\
线性细化、严格等价 & 同极化映射间的线性差；共同可逆输入与目标变换 & \hyperlink{section.2.3}{2.3}、\hyperlink{section.5.3}{5.3} \\
\(E_{0,z},E_{1,z},\gamma,\delta\) & 核心／余块的原始Walsh相位及其原点常数 & \hyperlink{section.5.4}{5.4}、\hyperlink{section.6.1}{6.1—6.3} \\
\(Q_*,H_+,\perp\)、中心积 & 固定八维核心、四维正平面、共同目标正交和及对应群拼接 & \hyperlink{section.1.5}{1.5}、\hyperlink{section.6.6}{6.6—6.8} \\
MM公式、EA等价、对偶闭合 & 既有求和工具及其适用／等价范围 & \hyperlink{section.8.3}{8.3—8.6} \\
\bottomrule
\end{longtable}


尤其不要混淆三个层次：\(D\)恒零是函数陈述；\(C=0\)是极化陈述；\(\sum(-1)^D=0\)是整数字符和陈述。也不要把8维核心的维数误当群阶的指数：核心群阶是\(2^{8+2}=1024\)。



\tableofcontents
\mainmatter


\chapter{从有限群到平方映射：读这篇论文需要的背景与全图}


本解析对应预印本 v0.1.0，日期为2026年10月17日。对应19页英文稿 \emph{Extremal Klein Four Counts in Semi-Extraspecial 2-Groups}；下文的定理、引理、公式编号沿用该稿，以便逐项核对。英文稿页码指这份19页稿；新版排版若页码移动，应以相同编号为准。本章对应英文稿第1节，尤其第1—3页。解析中的展开计算是对已有证明的教学解释，不自动变成论文新增结果，也不表示作者已经完成人工核验。



\Needspace{7\baselineskip}
\section{问题为什么不只是数元素}


有限群是带有结合乘法、有单位元且每个元素可逆的有限集合。本题首先数满足 \(g^2=1\) 的元素。这包括单位元本身，以及阶恰为2的元素。文中把这个集合记为



\[
I(G)=\{g\in G:g^2=1\},\qquad s(G)=|I(G)|.
\]



“有多少个这样的元素”并不能告诉我们“这些元素彼此怎样相乘”。两个不同的非单位自逆元素 \(g,h\)，若还满足 \(gh=hg\)，就组成



\[
\{1,g,h,gh\}\cong C_2\times C_2.
\]



这就是Klein四元群，也称四元子群。这里的 \(C_2\) 是二阶循环群。本论文的另一个计数是



\[
v_4(G)=|\{H\le G:H\cong C_2\times C_2\}|.
\]



它统计实际子群，不能把互为共轭的子群合并。论文研究：在一个结构受控的群类里，固定群阶并令 \(s(G)\) 达到最大，\(v_4(G)\) 还能怎样变化？达到最小 \(v_4\) 的群能否全部认出来？这个问题把“元素阶”与“交换关系”联系起来，因而适合用平方和交换子共同研究。



\Needspace{7\baselineskip}
\section{中心、导出子群、Frattini子群在这里各管什么}


\Needspace{5\baselineskip}
中心



\[
Z(G)=\{w\in G:wg=gw\text{ 对每个 }g\in G\}
\]



记录与所有元素交换的部分。交换子可取约定 \([g,h]=g^{-1}h^{-1}gh\)；导出子群 \(G'\) 是全部交换子生成的子群。把 \(G'\) 商掉，群变为交换群。因此 \(G'\) 衡量非交换性。



Frattini子群 \(\Phi(G)\) 是所有极大真子群的交。在有限2群中，常用公式是



\[
\Phi(G)=G^2G',
\]



其中 \(G^2\) 表示所有平方生成的子群。这个公式的含义是：商掉平方与交换子后，所得商群是一个初等交换2群，即若干个 \(C_2\) 的直积，也是一个 \(\mathbb F_2\) 向量空间。使用有限2群的极大子群都有指数2这一基础事实，每个极大子群都含平方与交换子，所以 \(G^2G'\le\Phi(G)\)。反过来，若 \(g\notin G^2G'\)，可在向量空间 \(G/(G^2G')\) 中选一个不含 \(g\) 像的超平面；其逆像是指数2的极大子群，且不含 \(g\)。于是 \(g\notin\Phi(G)\)，给出反向包含。这是英文命题2.1使用的标准有限2群事实。平方映射与special 2群的对应框架见[KK12, KK15]；本解析第2章另把本题所需的构造与同构对应完整展开，不要求读者先掌握那套理论。



非交换2群若满足 \(Z(G)=G'=\Phi(G)\)，且这个共同子群是初等交换群，称为special（特殊）2群。若共同子群进一步只有两个元素，称为extraspecial（额外特殊）2群。semi-extraspecial（半额外特殊）条件则要求：对中心中每个指数2子群 \(U\)，\(G/U\) 都是extraspecial。



\Needspace{5\baselineskip}
本论文同时施加



\[
Z(G)=G'=\Phi(G)=W\cong C_2^2,\qquad V=G/W\cong\mathbb F_2^{2n}.
\tag{英文(1.1)}
\]



所以 \(W\) 有4个元素；中心中有3个二阶子群，每个都是一个非零线性泛函 \(\lambda:W\to\mathbb F_2\) 的核。semi-extraspecial要求的恰好是这3个中央商都extraspecial。不能把“special且中心有4个元素”单独当作这个更强条件。



本题 \(|V|=2^{2n}\)，故



\[
|G|=|W||V|=2^{2n+2}.
\]



参数 \(n\) 是 \(V\) 的\textbf{半维数}。说“奇数分支”或“偶数分支”都是说 \(n\)，不是说 \(\dim_{\mathbb F_2}V\)；后者始终是偶数 \(2n\)。群阶也不是 \(2^{2n}\)，还须乘中心的4。



\Needspace{7\baselineskip}
\section{为什么商群是向量空间，平方却取值于中心}


\(V=G/W\) 的一个元素是陪集 \(gW\)，不是选定的群元素。由于 \(G^2\le\Phi(G)=W\)，每个陪集平方都是单位陪集；又因 \(G'\le W\)，这个商群交换。于是把陪集乘法写成加法，就成为二元向量空间：\(gW+hW\) 表示 \(ghW\)。



群中平方 \(g^2\) 落在 \(W\)。换同一陪集的代表 \(g\) 为 \(gw\)，利用 \(w\) 中心且 \(w^2=1\)，有



\[
(gw)^2=g^2w^2=g^2.
\]



\Needspace{5\baselineskip}
因此可以定义



\[
Q:V\longrightarrow W,\qquad Q(gW)=g^2.
\]



把 \(W\) 也写成加法向量空间，单位元就写作 \(0\)。\(Q(v)=0\) 等价于该陪集中的4个元素全部自逆。



\Needspace{5\baselineskip}
平方不是线性函数。它的“不线性部分”是



\[
B(x,y)=Q(x+y)+Q(x)+Q(y).
\]



在本群类中，这正是交换子映射。\(B\) 是双线性且 \(B(x,x)=0\)，这种双线性映射称为交替的。因特征为2，交替形式的矩阵也是对称的，但对角为0；不要把“交替”误解为矩阵只有零。



满足 \(Q(0)=0\)、且这个 \(B\) 双线性的函数，称为 \(\mathbb F_2\) 二次映射。它允许一次项。比如 \(xy+x\) 与 \(xy\) 有相同极化，但函数不同。在二元坐标上 \(x^2=x\)，所以“对角平方项”也能表现为一次函数。后面必须保存的线性细化，指的就是这些同极化而不同平方映射的信息。



\Needspace{7\baselineskip}
\section{三个标量分量与非退化条件}


取 \(W\) 的一个二元基，把 \(Q\) 写成 \((q_1,q_2)\)。由于 \(W^*\) 恰有3个非零线性泛函，全部非零标量分量为



\[
q_1,\qquad q_2,\qquad q_3=q_1+q_2.
\]



它们的极化分别记为 \(B_1,B_2,B_3=B_1+B_2\)。标量双线性形式的根空间（radical）是



\[
\operatorname{rad}(B_i)=\{x:B_i(x,y)=0\text{ 对所有 }y\}.
\]



“非退化”就是根空间为零，等价于其矩阵可逆。在中央商 \(G/\ker\lambda\) 中，一个陪集在中心里，当且仅当它的交换子被 \(\lambda\) 全部杀掉。因此semi-extraspecial恰好给出：\(B_1,B_2,B_3\) 全部非退化。这个翻译是整个方法的入口，不只是为了选择方便记号。



非退化交替形式只能存在于偶维空间。可由以下过程看出：选非零 \(u\)，非退化性给 \(v\) 使 \(B(u,v)=1\)；它们张成一个非退化二维平面，再对其正交补重复。得到的基叫辛基。这个二维平面分解会在字符和、Arf符号和矩阵笔论证中反复出现。



\Needspace{7\baselineskip}
\section{两个主定理的准确内容}


\Needspace{5\baselineskip}
设



\[
h=2^n,\qquad N=h^2=2^{2n},\qquad
M_n=\frac{h^4+9h^3+36h^2-36h+8}{24}.
\tag{英文(1.2)}
\]



英文定理1.1（基线第2页）先说：\(n=1\) 没有满足上述条件的群；每个 \(n\ge2\) 都有这样的群，且



\[
s(G)\le N+3h.
\]



这个上界在每个允许的 \(n\) 都能达到。然后只在\textbf{最大自逆元素数分支} \(s(G)=N+3h\) 上，给四元子群数的最优界：



\begin{longtable}{@{}p{0.3033\linewidth}@{\hspace{.025\linewidth}}p{0.3033\linewidth}@{\hspace{.025\linewidth}}p{0.3033\linewidth}@{}}
\toprule
半维数 & 下界 & 上界 \\
\midrule
\endhead
\(n=2\) & \(M_n\) & \(M_n+2^{3n-2}\) \\
奇数 \(n\ge3\) & \(M_n-2^{3n-5}\) & \(M_n+2^{3n-5}\) \\
偶数 \(n\ge4\) & \(M_n-2^{3n-4}\) & \(M_n+2^{3n-2}\) \\
\bottomrule
\end{longtable}


每一行的两个端点都能达到，但没有声称端点间每个整数都能达到。\(n=2\) 更精确，只有53与69这两个值。最初三个实例为



\begin{longtable}{@{}p{0.1720\linewidth}@{\hspace{.025\linewidth}}p{0.1720\linewidth}@{\hspace{.025\linewidth}}p{0.1720\linewidth}@{\hspace{.025\linewidth}}p{0.1720\linewidth}@{\hspace{.025\linewidth}}p{0.1720\linewidth}@{}}
\toprule
\(n\) & \(|G|\) & \(s_{\max}\) & 最小 \(v_4\) & 最大 \(v_4\) \\
\midrule
\endhead
2 & 64 & 28 & 53 & 69 \\
3 & 256 & 88 & 431 & 463 \\
4 & 1024 & 304 & 4371 & 5651 \\
\bottomrule
\end{longtable}


定理1.2（基线第2页；完整证明在第14页）不分类全部群，而只分类\textbf{偶数 \(n\ge4\)、最大 \(s\)、且达到下界}的群。令



\[
K=\mathbb F_4=\mathbb F_2[\alpha],\qquad \alpha^2+\alpha+1=0,
\]



且 \(\operatorname{Tr}:K\to\mathbb F_2\) 为 \(u\mapsto u+u^2\)。把迹值经 \(\mathbb F_2\subset K\) 看成 \(K\) 值，定义



\[
Q_*(a,b,c,d)=ad+bc+\operatorname{Tr}(ac)+a+\alpha c,
\qquad H_+(u,v)=uv.\tag{英文(1.5)}
\]



\(Q_*\) 的定义域是 \(K^4\)，即8维二元空间；\(H_+\) 的定义域是 \(K^2\)，即4维二元空间。定理说，下列三件事等价：



\begin{enumerate}


\item \(s=N+3h\) 且 \(v_4=M_n-2^{3n-4}\)。

\item 完整平方映射与 \(Q_*\perp H_+^{\perp r}\) 严格等价，其中 \(r=(n-4)/2\)；全部块取值于\textbf{同一个} \(W=K\)。

\item 群是 \(Q_*\) 对应群与 \(r\) 个64阶Heisenberg群 \(H(K)\) 沿共同中心 \(W\) 的中心积。

\end{enumerate}


严格等价只允许可逆二元线性映射 \(S,T\)，满足



\[
\widetilde Q(Sv)=TQ(v).\tag{英文(1.4)}
\]



它不允许随意添加线性映射、平移原点、或为不同块各取不相容的目标换基。定理断言每个固定偶数 \(n\ge4\) 的\textbf{整个群}只有一个同构类型；并不声称嵌入的八维核心唯一，或中心分解唯一。维数确实吻合：\(8+4r=2n\)。\(n\) 为奇数时这个表达式本身就不适配；论文还给出更强的非分解例，不能把偶数结论直接外推。



\Needspace{7\baselineskip}
\section{有限域坐标不改变二次性的基域}


\(K\) 有四个元素 \(0,1,\alpha,\alpha^2\)，其中 \(\alpha^2=\alpha+1\)，\(\alpha^3=1\)。用它写坐标能把两个二元坐标压缩为一个，但不会自动使每个 \(\mathbb F_2\) 二次映射成为 \(K\) 上齐次二次型。



\(K\) 齐次二次型须满足 \(q(\lambda x)=\lambda^2q(x)\)，且极化对 \(K\) 双线性。\(Q_*\) 中普通一次项 \(a+\alpha c\)、以及嵌入 \(K\) 的 \(\operatorname{Tr}(ac)\)，不能未经检查就按 \(K\) 齐次二次型处理。对 \(\mathbb F_2\) 而言，它们则分别是允许的一次项与二次项。



有限域上的迹配对 \(\operatorname{Tr}(xy)\) 是非退化二元配对。选定它后，频率变量可写成有限域元素；其本质仍是二元对偶空间里的线性泛函。后文从二元矩阵转到 \(K\) 坐标，都必须同时运输这个配对，不能对三个分量各自偷偷换成不同配对。



\Needspace{7\baselineskip}
\section{全文证明路线和符号导航}


第2章先证明群与完整平方映射的对应，再把 \(v_4\) 化成零点集上有序对的计数。第3章把这个计数写成三阶Walsh矩；再用三对偶和 \(D\) 的根空间与符号得到界及等号条件。第4章构造每个端点，完成定理1.1。后续章节针对偶数下界：先由等号条件提取八维极化块，再筛选完整线性细化、归一96个带坐标映射，最后处理核心与整个余块的相位，证明定理1.2。范围反例与附录说明哪些更强推断并不成立。



\Needspace{5\baselineskip}
反复出现的记号如下：



\begin{longtable}{@{}p{0.3033\linewidth}@{\hspace{.025\linewidth}}p{0.3033\linewidth}@{\hspace{.025\linewidth}}p{0.3033\linewidth}@{}}
\toprule
记号 & 所在对象与含义 & 首次关键位置 \\
\midrule
\endhead
\(W\) & 中心、导出及Frattini子群；二元维数2 & 英文(1.1) \\
\(V\) & \(G/W\)；二元维数 \(2n\) & 英文(1.1) \\
\(Q,B\) & 完整平方映射及其极化（交换子映射） & 命题2.1 \\
\(S,z\) & \(Q^{-1}(0)\) 及其大小 & 引理2.3之前 \\
\(T_0\) & \(x,y,x+y\) 同在 \(S\) 的有序对数 & 引理2.3 \\
\(L\) & \(Q\) 恒零的二维二元子空间数 & 英文(2.3) \\
\(J_2\) & 群中相互交换的有序自逆元素对数 & 引理2.3 \\
\(\mathcal W_i(\xi)\) & 第 \(i\) 个标量分量的Walsh值 & 第3节开头 \\
\(\mathcal P\) & \(\sum_\xi\mathcal W_1\mathcal W_2\mathcal W_3\) & 命题3.2 \\
\(q_i^*\) & 最大 \(s\) 分支的实际Walsh对偶 & 英文(3.3) \\
\(D\) & 三个实际对偶的二元函数和 & 英文(3.4) \\
\(C,\rho\) & \(D\) 的极化及根空间维数 & 英文(3.4) \\
\(A,p\) & \(B_1^{-1}B_2\) 及 \(p(t)=t^2+t+1\) & 命题3.4 \\
\(Q_*,H_+\) & 八维核心、四维正端块 & 英文(1.5) \\
\bottomrule
\end{longtable}


最后要分清论证与文献定位。平方映射模型与中心积使用[KK12, KK15]的框架，交替矩阵笔与局部块属于[LM18, PFMR14, MPG13]涉及的已有理论；标量Walsh对偶及一般Maiorana–McFarland公式使用[AFH24, CMP18]。给定等号极化后96个细化和符号筛选，是已有公式与经典判据的短推导，不是新的通用对偶理论。论文组织的是这个指定群类的最优界、准确等号及偶数分支的全体严格同构分类。已核查文献尚未直接替代这个完整合成结论，并不等于原创性已被绝对确认。AFH与CMP全文已经取得并完成对照，不是当前获取缺口。



\chapter{平方映射、群同构和每一个计数倍率}


本章逐步解释英文第2节：命题2.1在基线第3—4页，引理2.2与2.3在第4页。它负责两个转换：先把有限群问题准确转换成二次映射问题，再把实际四元子群数转换成零点集计数。后续字符和方法只对转换后的问题运算；这里的倍率若错了，后面的数值全部会错。



\Needspace{7\baselineskip}
\section{从群出发：平方为何良定义，极化为何等于交换子}


由于 \(G'=W\le Z(G)\)，群的交换子都在中心，这称为幂零类至多2。又因 \(G/W\) 是初等交换2群，\(g^2\in W\)；而 \(W\) 的每个元素平方为单位元，所以 \(g^4=1\)。这个群指数整除4。



平方映射定义为 \(Q(gW)=g^2\)。若换代表为 \(gw\)，有



\[
(gw)^2=g^2w^2=g^2,
\]



第一步使用 \(w\in Z(G)\)，第二步使用 \(w^2=1\)。只知道“\(w\) 在中心”还不够。类似地，换中心倍数不会改变交换子。



\Needspace{5\baselineskip}
类2群中的平方恒等式是



\[
(gh)^2=g^2h^2[h,g].
\]



由于这里交换子阶至多2，\([h,g]=[g,h]^{-1}=[g,h]\)。把中心写成加法，得到



\[
Q(gW+hW)+Q(gW)+Q(hW)=[g,h].
\]



因交换子中心化，\([g_1g_2,h]=[g_1,h][g_2,h]\)，另一变量也如此，所以 \(B\) 双线性。\([g,g]=1\) 给 \(B(x,x)=0\)。这同时解释了“二次映射”和“交替极化”并非额外猜测，而是群乘法直接带来的结构。



取 \(\lambda\in W^*\setminus\{0\}\)，\(U=\ker\lambda\)。在 \(G/U\) 中，\(gU\) 与全部元素交换，当且仅当全部 \([g,h]\) 属于 \(U\)，即



\[
\lambda B(gW,hW)=0\quad\text{对每个 }h.
\]



所以 \((G/U)\) 的中心超出 \(W/U\) 的部分，正好由 \(\operatorname{rad}(\lambda B)\) 描述。由于 \((G/U)'=W/U\) 已是二阶，而平方也落在 \(W/U\)，Frattini公式给 \(\Phi(G/U)=W/U\)。因此 \(G/U\) 为extraspecial当且仅当这个根空间为零。于是semi-extraspecial条件与三个标量极化可逆完全对应。



\Needspace{7\baselineskip}
\section{反过来：给定二次映射怎样造出一个群}


英文命题2.1的逆向保证后面构造的函数确实来自满足假设的群。设 \(V\) 是非零二元空间，\(W=\mathbb F_2^2\)，\(Q:V\to W\) 二次，三个非零标量极化均非退化。



\Needspace{5\baselineskip}
选坐标后，任何这样的函数可写为



\[
Q(v)=\sum_i a_i v_i+\sum_{i<j}c_{ij}v_iv_j,
\]



其中系数 \(a_i,c_{ij}\in W\)。没有重复写 \(v_i^2\)，因为在二元输入上 \(v_i^2=v_i\)。定义双线性映射



\[
F(v,u)=\sum_i a_iv_iu_i+\sum_{i<j}c_{ij}v_i u_j.
\]



它一般不对称，但 \(F(v,v)=Q(v)\)。在集合 \(V\times W\) 上规定



\[
(v,w)(u,t)=(v+u,w+t+F(v,u)).\tag{英文(2.1)}
\]



结合律的关键计算只有一个：对第三个向量 \(r\)，两种括号中的额外中心项分别为



\[
F(v,u)+F(v+u,r),\qquad F(u,r)+F(v,u+r).
\]



双线性把它们都展开为 \(F(v,u)+F(v,r)+F(u,r)\)，所以相等。\((0,0)\) 是单位元，\((v,w+Q(v))\) 是 \((v,w)\) 的逆元。平方与交换子分别为



\[
(v,w)^2=(0,Q(v)),\qquad
[(v,w),(u,t)]=(0,F(v,u)+F(u,v))=(0,B(v,u)).
\]



可见整个群阶是 \(4\cdot2^{\dim V}\)，平方也确实是预定的 \(Q\)。



还要检查三个结构子群。若 \((v,w)\) 中心化，则 \(B(v,u)=0\) 对全部 \(u\)；取任意非零 \(\lambda\)，非退化性迫 \(v=0\)。故中心恰好是 \(\{0\}\times W\)。若所有 \(B(v,u)\) 张成 \(W\) 的真子空间，会有非零 \(\lambda\) 把它们全部杀掉，导致 \(\lambda B=0\)，违背非退化。因此导出子群是整个 \(W\)。最后平方已在 \(W\) 内，\(\Phi(G)=G^2G'=W\)。上一节中央商的论证又给semi-extraspecial。



这里“每个非零标量极化非退化”同时用于确定中心、保证导出子群不缩小以及核对全部中央商；不能只在构造时查其中一个矩阵。



\Needspace{7\baselineskip}
\section{为什么完整平方映射严格等价，恰好对应群同构}


群同构必保中心，故诱导可逆线性映射 \(S:V\to\widetilde V\)、\(T:W\to\widetilde W\)。又因同构保平方，必有



\[
\widetilde Q(Sv)=TQ(v).
\]



反方向不能仅凭一句“显然”略过：平方和交换子是否已经记录全部群关系？选 \(V\) 的基 \(e_1,\ldots,e_{2n}\)，在群里选提升 \(g_i\)。再选 \(W\) 的基。所有关系可写成中心元素阶2、中心化、\(g_i^2=Q(e_i)\)、\([g_i,g_j]=B(e_i,e_j)\)。



任意字可用交换关系把 \(g_i\) 按顺序排列，再用平方关系把每个指数减到0或1。因此所有元素都能写成



\[
w g_1^{a_1}\cdots g_{2n}^{a_{2n}},\qquad w\in W,\quad a_i\in\{0,1\}.
\]



这样最多有 \(4\cdot2^{2n}\) 个字，而实际群恰好这么大，故没有遗漏会造成额外群类型的关系。若给定 \(S,T\) 满足严格等价，把 \(g_i\) 送到 \(Se_i\) 的任意提升，把中心按 \(T\) 送过去，就保全部上述关系。得到等阶群间的满同态，因而是同构。



这个证明说明了为什么只分类 \(B\) 不够。若在 \(Q\) 上添加线性映射 \(L\)，其极化不变，但生成元平方可能变了。第4章的 \(uv\) 与 \(uv+u\) 有相同 \(B\)、相同最大 \(s\)，却有不同 \(v_4\)，所以绝不可能同构。它们正好是“交换子数据并不决定完整群”的具体例子。



\Needspace{7\baselineskip}
\section{正交和使用共同目标，群里对应中心积}


若 \(V=V_0\oplus V_1\) 且 \(B(V_0,V_1)=0\)，那么对 \(x\in V_0,y\in V_1\)，极化定义直接给



\[
Q(x+y)=Q(x)+Q(y).
\]



这写成 \(Q=Q_0\perp Q_1\)。两个值都在同一个 \(W\) 中相加，而不是组成 \(W\oplus W\) 的一对。



在相应群里，\(V_0,V_1\) 的逆像 \(G_0,G_1\) 彼此交换、交于 \(W\)、共同生成 \(G\)。乘法映射 \(G_0\times G_1\to G\) 的核是



\[
\{(w,w):w\in W\}
\]



（一般写成 \((w,w^{-1})\)；这里 \(w^{-1}=w\)）。所以是中心积，而不是直积。对 \(r+1\) 块，核是



\[
\mathcal N=\{(w_0,\ldots,w_r)\in W^{r+1}:\textstyle\sum_iw_i=0\}.
\]



直接检查阶数也能看出区别：两块若商维数分别为 \(d_0,d_1\)，直积阶数是 \(2^{d_0+d_1+4}\)，再除中心核的4得到 \(2^{d_0+d_1+2}\)，这才与单个共同中心对应的平方映射一致。英文引理2.2给出完整证明，并对应[KK12, Lemma 2.4；所指2012年v1 PDF第6页]。不能把另一个呈现版本的编号混进来。



\Needspace{7\baselineskip}
\section{零点集到群元素：\texorpdfstring{\(s=4z\)}{s=4z} 与 \texorpdfstring{\(J_2=16T_0\)}{J2=16T0}}


\Needspace{5\baselineskip}
记



\[
S=Q^{-1}(0),\quad z=|S|,
\quad T_0=|\{(x,y)\in S^2:x+y\in S\}|.
\]



\(S\) 一般不是子空间；后面的 \(T_0\) 正是在数它还保留多少加法关系。群里定义



\[
J_2=|\{(g,h)\in I(G)^2:gh=hg\}|.
\]



每个 \(v\in S\) 的陪集里4个元素全部平方为1。\(v\notin S\) 的陪集没有这样的元素。因此 \(s=4z\)。特别地，\(V\) 的零向量对应的是整个中心 \(W\)，中心的单位元和3个非单位元素都已经计在这一个零向量的4个提升里。



现在固定有序向量对 \((x,y)\in S^2\)，各选群提升 \(g,h\)。它们自逆且



\[
[g,h]=B(x,y)=Q(x+y),
\]



因为 \(Q(x)=Q(y)=0\)。交换子与选择哪个中心提升无关。因此若 \(x+y\in S\)，4种 \(g\) 与4种 \(h\) 的全部16种组合都交换；否则没有一种交换。这证明



\[
s=4z,\qquad J_2=16T_0.
\]



16是两个\textbf{有序}变量各自有4种提升，并不是只有乘积有16种选择。这里无需固定一个提升后再猜另一个：两个选择是独立的。



\Needspace{7\baselineskip}
\section{从交换有序对到四元子群：\texorpdfstring{\(3s-2\)}{3s-2} 与6的来源}


在 \(J_2\) 统计的有序对中，先剥离不能生成四元群的退化情况：



\begin{itemize}


\item 对角对 \((g,g)\)，共有 \(s\) 个，包括 \((1,1)\)。

\item 第一项是单位元、第二项非单位元，共 \(s-1\) 个。

\item 第二项是单位元、第一项非单位元，共 \(s-1\) 个。

\end{itemize}


这三类不交，总数 \(3s-2\)。除此之外，\(g,h\) 是不同的非单位自逆元素，且交换，因此生成唯一四元群。反过来，一个四元群有3个非单位元，任意挑不同的两个并安排顺序，有 \(3\cdot2=6\) 对，全部生成它。于是



\[
J_2=3s-2+6v_4(G).\tag{英文(2.2)的一部分}
\]



“6”不是共轭计数，也不是选无序两元素的3。若改成无序对，整个公式必须一起改。本论文一直用有序对。



在向量空间里照样数：\(T_0\) 的退化对是 \(x=y\)、\(x=0\ne y\)、\(y=0\ne x\)，合计 \(3z-2\)。剩余 \(x,y\) 是线性无关的两个二元向量；它们的张成空间恰有 \(0,x,y,x+y\)，全部在 \(S\)。记这种二维全零子空间数为 \(L\)，则每个又给6个有序基，所以



\[
T_0=3z-2+6L.
\]



\Needspace{5\baselineskip}
代入两边的提升公式，得到



\[
v_4(G)=\frac{16T_0-3s+2}{6}=16L+6z-5.
\tag{英文(2.3)}
\]



英文引理2.3第4页同时给这几式。它们解释了后文为何只需计算 \(T_0\)，不必直接枚举群的全部子群。



\Needspace{7\baselineskip}
\section{教学补充：直接按中心交分类，重新看 \texorpdfstring{\(16L+6z-5\)}{16L+6z-5}}


这不是新增论文结论，而是同一个恒等式的另一种核对，也能提醒我们“零平面”与“群中四元子群”不是一一对应。



设 \(H\cong C_2^2\)。按 \(H\cap W\) 的维数分类：



\begin{enumerate}


\item 若维数2，则 \(H=W\)，只有1个。

\item 若维数1，\(H\) 在 \(V\) 中的像是一条非零零线 \(\{0,x\}\)，有 \(z-1\) 种。先选 \(H\cap W\)，中心里有3条线。固定这条线 \(K\) 后，\(x\) 的4个提升按 \(K\) 分成2个陪集，每个选择给一个 \(H=K\cup gK\)。所以总数 \(6(z-1)\)。

\item 若维数0，像是一个二维全零子空间 \(U\)。它的群逆像是一个四维初等交换群：平方全为1，且极化在 \(U\) 上为零。\(H\) 是中心 \(W\) 在其中的补空间。选一个补空间后，其余补空间是线性映射 \(U\to W\) 的图像；这类映射共 \(2^{2\cdot2}=16\) 个。故每个 \(U\) 提供16个 \(H\)，合计 \(16L\)。

\end{enumerate}


三者相加是 \(1+6(z-1)+16L\)。这里的16是“每个零平面在群中有16个补空间提升”，而前一节的16是“每个有序零对有16个提升”；它们来自相容的线性计数，但不能不说明对象就互相替换。



至此，群论侧的工作完成。接下来真正未知的是零点集里满足 \(x+y\in S\) 的有序对数 \(T_0\)。它涉及三个零点条件，正好会变成一个三阶Fourier／Walsh矩。



\chapter{三阶Walsh矩、根空间亏损与最优界}


本章解释英文第3节，基线第5—7页。它包括引理3.1、命题3.2、引理3.3、命题3.4及3.5，给出数值主定理的整个必要性部分。这里的关键不是一次复杂求和，而是两层转换：先把零点集加法关系化成三个Walsh变换的乘积，再把这个乘积的总和化成一个二次函数 \(D\) 的字符和。最后，交替矩阵笔限制 \(D\) 可以有多大的根空间。



\Needspace{7\baselineskip}
\section{字符正交性是离散版的“筛选器”}


对 \(\xi\in V^*\)，函数 \(x\mapsto(-1)^{\xi(x)}\) 叫一个二元加法字符。若 \(x=0\)，所有字符都取1；若 \(x\ne0\)，可选线性泛函 \(\eta\) 使 \(\eta(x)=1\)，把 \(\xi\) 与 \(\xi+\eta\) 配对，符号就相反。因此



\[
\sum_{\xi\in V^*}(-1)^{\xi(x)}=
\begin{cases}N,&x=0,\\0,&x\ne0.\end{cases}
\]



同样，对一个非零 \(\xi\)，\(\sum_x(-1)^{\xi(x)}=0\)。Fourier变换取约定



\[
\widehat f(\xi)=\sum_{x\in V}f(x)(-1)^{\xi(x)}.
\]



二次函数 \(q_i\) 的Walsh变换就是 \((-1)^{q_i}\) 的Fourier变换：



\[
\mathcal W_i(\xi)=\sum_x(-1)^{q_i(x)+\xi(x)}.
\]



零频 \(\mathcal W_i(0)\) 又叫原始Gauss和。非零频率等于在 \(q_i\) 上加一个线性项后再做Gauss和。因此研究全部Walsh值，就是同时研究全部线性细化的符号。



\Needspace{7\baselineskip}
\section{引理3.1：二次字符和为何只有零或特定正负大小}


设 \(d(0)=0\)，其极化根空间是 \(R\)，\(\dim R=\rho\)。在 \(R\) 上，极化为零，所以



\[
d(r+r')=d(r)+d(r'),
\]



即 \(d|_R\) 为线性函数。如果有 \(r\in R\) 满足 \(d(r)=1\)，则



\[
d(x+r)=d(x)+d(r)+B_d(x,r)=d(x)+1.
\]



平移 \(x\mapsto x+r\) 给求和项两两抵消，所以 \(\sum_x(-1)^{d(x)}=0\)。注意：这是\textbf{整数和为零}，并不是 \(d\) 函数恒为零。



若 \(d|_R=0\)，则 \(d(x+r)=d(x)\)，所以它在商空间 \(V/R\) 上良定义。商极化非退化。设这个商维数为 \(2m=2n-\rho\)，在辛基中每个二维块的函数形如



\[
xy+ax+by,\qquad a,b\in\mathbb F_2.
\]



直接把四个输入相加，或先对 \(y\) 求和，可得



\[
\sum_{x,y\in\mathbb F_2}(-1)^{xy+ax+by}=2(-1)^{ab}.
\]



正交块的和是各块字符和的乘积，因此非退化 \(2m\) 维函数的和为 \(\pm2^m\)。从商提升时，每个向量有 \(2^\rho\) 个原像，所以



\[
\sum_x(-1)^{d(x)}=\pm2^\rho2^{(2n-\rho)/2}
=\pm2^{n+\rho/2}.
\]



符号等于辛平面中 \(\sum ab\) 的奇偶；这就是二元非退化二次型的Arf符号。这里不必先背Arf分类：这个四点计算已足以解释本章需要的求和判据。由于改变基不改变整数和，符号也不依赖所选辛基。



对非退化 \(q_i\)，根空间维数0，任意线性细化仍有同一个非退化极化，故每个Walsh值大小均为 \(h=2^n\)，从而



\[
\mathcal W_i(\xi)^2=N.
\]



“bent”在本二元二次背景里正是全部Walsh值具有这个恒定绝对值。它不表示原始Gauss和必为正；符号仍需另外处理。



\Needspace{7\baselineskip}
\section{最大自逆元素数分支从哪里来}


一个向量同时满足 \(q_1(x)=q_2(x)=0\) 当且仅当 \(Q(x)=0\)。利用二元零值指示函数，有



\[
1_S(x)=\frac{(1+(-1)^{q_1(x)})(1+(-1)^{q_2(x)})}{4}
=\frac{1+(-1)^{q_1(x)}+(-1)^{q_2(x)}+(-1)^{q_3(x)}}4.
\]



令 \(a_0=\mathcal W_1(0)+\mathcal W_2(0)+\mathcal W_3(0)\)，求和得



\[
z=\frac{N+a_0}{4},\qquad s=4z=N+a_0.
\]



每个原始和是 \(\pm h\)，故 \(s\le N+3h\)，且达到这个值当且仅当三个原始和\textbf{各自都为 \(+h\)}。这是英文命题3.2前半。当前只证明了上界及其判据，尚未证明每个 \(n\) 都可达；可达性留到第4章。



这种全局正号并不自动传到任意正交分块。一个标量全局Gauss和是两个块Gauss和的乘积，两个负号也能乘成正号。后面处理核心与余块时必须重新利用下界等号，而不能在这里先把各块的符号都设为正。



\Needspace{7\baselineskip}
\section{完整展开三阶Walsh矩}


令 \(f=1_S\)。字符正交性给



\[
\begin{aligned}
\sum_\xi\widehat f(\xi)^3
&=\sum_{x,y,z}f(x)f(y)f(z)
  \sum_\xi(-1)^{\xi(x+y+z)}\\
&=N\sum_{x,y}f(x)f(y)f(x+y)=NT_0.
\end{aligned}
\]



这就是为何要用“三阶”：\(x,y,x+y\) 三个零点条件产生三个Fourier因子。设



\[
U(\xi)=\mathcal W_1(\xi)+\mathcal W_2(\xi)+\mathcal W_3(\xi),\qquad
\mathcal P=\sum_\xi\mathcal W_1(\xi)\mathcal W_2(\xi)\mathcal W_3(\xi).
\]



\Needspace{5\baselineskip}
指示函数公式变换后是



\[
\widehat f(\xi)=\frac{N1_{\xi=0}+U(\xi)}4.
\]



这里常数函数1的Fourier变换只在零频有值 \(N\)。现在逐项展开立方，不预先假设最大 \(s\)：



\[
\sum_\xi(N1_{\xi=0}+U(\xi))^3
=N^3+3N^2a_0+3Na_0^2+\sum_\xi U(\xi)^3.
\]



对任意固定频率，暂记 \(w_i=\mathcal W_i(\xi)\)。因 \(w_i^2=N\)，



\[
\begin{aligned}
(w_1+w_2+w_3)^3
&=\sum_iw_i^3+3\sum_{i\ne j}w_i^2w_j+6w_1w_2w_3\\
&=N\sum_iw_i+3N\bigl(2\sum_jw_j\bigr)+6w_1w_2w_3\\
&=7NU(\xi)+6w_1w_2w_3.
\end{aligned}
\]



系数7来自 \(1+3\cdot2\)，不能把中间的有序下标和误当3项。再由 \(q_i(0)=0\) 和字符正交性，



\[
\sum_\xi\mathcal W_i(\xi)=N(-1)^{q_i(0)}=N,
\qquad \sum_\xi U(\xi)=3N.
\]



所以 \(\sum U^3=21N^2+6\mathcal P\)。除以指示函数的 \(4^3=64\)，再除以 \(N\)，得到



\[
T_0=\frac{N^2+3Na_0+3a_0^2+21N}{64}
       +\frac{3\mathcal P}{32N}.
\tag{英文(3.1)}
\]



在最大 \(s\) 分支上 \(a_0=3h\)。因为 \(N=h^2\)，上式化成



\[
T_0=\frac{h^4+9h^3+48h^2}{64}+\frac{3\mathcal P}{32N}.
\]



第2章给 \(v_4=(16T_0-3s+2)/6\)。代入 \(s=h^2+3h\) 后，分子中 \(48h^2/4-3h^2=9h^2\)，于是



\[
v_4(G)=\frac{h^4+9h^3+36h^2-36h+8}{24}
         +\frac{\mathcal P}{4N}
=M_n+\frac{\mathcal P}{4N}.
\tag{英文(3.2)}
\]



到这里，所有群依赖都压缩进一个整数 \(\mathcal P\)。这一步没有丢失任何线性项，也没有假设群同构类型已经知道。



\Needspace{7\baselineskip}
\section{从三个Walsh乘积到一个二次函数 \texorpdfstring{\(D\)}{D}}


选二元坐标，把线性泛函 \(\xi\) 写成列向量，使 \(\xi(x)=\xi^{\mathsf T}x\)。\(B_i\) 是 \(q_i\) 极化矩阵。对 \(t=B_i^{-1}\xi\)，



\[
q_i(x+t)=q_i(x)+q_i(t)+\xi(x).
\]



\Needspace{5\baselineskip}
平移求和变量给



\[
\mathcal W_i(\xi)=\mathcal W_i(0)(-1)^{q_i(B_i^{-1}\xi)}.
\]



一般若 \(\mathcal W_i(0)=h(-1)^{\epsilon_i}\)，\textbf{实际Walsh对偶}为 \(q_i(B_i^{-1}\xi)+\epsilon_i\)，原点相位不能删。最大 \(s\) 分支恰好使全部 \(\epsilon_i=0\)，此时才可记



\[
q_i^*(\xi)=q_i(B_i^{-1}\xi),\quad
\mathcal W_i(\xi)=h(-1)^{q_i^*(\xi)},\quad q_i^*(0)=0.
\tag{英文(3.3)}
\]



[AFH24, Theorem 2.1，印刷页509—510]从二次型的型别刻画这一标量对偶，包括二元线性细化；上面是本文实际使用的配平方计算。在基域 \(\mathbb F_2\) 上，一次项不违反其二次型函数框架。不能反过来因坐标用了 \(\mathbb F_4\)，就把尚未齐次化的向量映射直接套用较大基域的齐次二次型结论。



\Needspace{5\baselineskip}
定义



\[
D=q_1^*+q_2^*+q_3^*,\qquad
C=B_1^{-1}+B_2^{-1}+(B_1+B_2)^{-1},\qquad
\rho=\dim\ker C.
\tag{英文(3.4)}
\]



\(D\) 的加法在 \(\mathbb F_2\) 中逐点进行。每个 \(q_i^*\) 的极化为 \(B_i^{-1}\)，因为代入 \(B_i^{-1}\) 后矩阵成为 \((B_i^{-1})^{\mathsf T}B_iB_i^{-1}=B_i^{-1}\)。它仍交替：令 \(\xi=B_ix\)，就有 \(\xi^{\mathsf T}B_i^{-1}\xi=x^{\mathsf T}B_ix=0\)。所以 \(C\) 是 \(D\) 的交替极化。



三个Walsh值的乘积变成 \(h^3(-1)^{D(\xi)}\)，故



\[
\mathcal P=h^3\mathcal G(D),\qquad
\mathcal G(D)=\sum_\xi(-1)^{D(\xi)}.\tag{英文(3.5)}
\]



必须始终区分三件事：\(D\equiv0\) 是函数恒零；\(C=0\) 只说 \(D\) 线性（这里 \(D(0)=0\)）；\(\mathcal G(D)=0\) 是正负项数量一样。特别地，\(D\equiv0\) 时 \(\mathcal G(D)=N\)，绝非0。\(C=0\) 且 \(D\) 为非零线性函数时才有 \(\mathcal G(D)=0\)。



引理3.1现可重用：若 \(D\) 在 \(\ker C\) 上不恒零，\(\mathcal P=0\)；否则



\[
\mathcal P=\pm2^{3n}\,2^{n+\rho/2}=\pm2^{4n+\rho/2}.
\]



所以最优界取决于：\(\rho\) 最大可以多大，以及这个最大根维数能否给负号。



\Needspace{7\baselineskip}
\section{为什么要研究 \texorpdfstring{\(A=B_1^{-1}B_2\)}{A=B1-1B2}}


令 \(A=B_1^{-1}B_2\)，则 \(B_2=B_1A\)、\(B_3=B_1(I+A)\)。三个极化可逆保证 \(A\) 与 \(I+A\) 可逆。这排除特征多项式的线性因子 \(t\) 与 \(t+1\)。



置 \(p(t)=t^2+t+1\)。矩阵计算给



\[
\begin{aligned}
C&=[I+A^{-1}+(I+A)^{-1}]B_1^{-1}\\
&=[A(I+A)]^{-1}(A^2+A+I)B_1^{-1}.
\end{aligned}
\]



第二行可把第一行左乘 \(A(I+A)\) 核对：分子分别为 \(A+A^2\)、\(I+A\)、\(A\)，相加是 \(A^2+A+I\)。前面的因子可逆，所以



\[
C=[A(I+A)]^{-1}p(A)B_1^{-1},\qquad
\ker C=B_1\ker p(A).\tag{英文(3.6)}
\]



注意根空间是位于频率空间的 \(B_1\ker p(A)\)，不能把 \(B_1\) 静默删掉；只有维数相等是立即成立的。



为了限制 \(\ker p(A)\)，把 \(V\) 看成 \(\mathbb F_2[t]\) 模：多项式 \(f(t)\) 通过 \(f(A)\) 作用。“初等因子 \(f^k\) 的循环块”指一个同构于 \(\mathbb F_2[t]/(f^k)\) 的子模，二元维数是 \(k\deg f\)。\(f\) 是不可约多项式。重复简单 \(p\) 块与单个 \(p^2\) 块不是同一事：前者 \(p(A)=0\)，后者 \(p(A)\ne0\) 但平方为零，尽管特征多项式可以相同。



这里的矩阵笔（pencil）指 \(B_1,B_2\) 及它们的线性组合。用 \(A\) 把它改写成模，是已有交替矩阵笔理论的语言。为了不把关键限制藏在引用后，英文引理3.3给了下述成对证明。



\Needspace{7\baselineskip}
\section{引理3.3：为什么每个初等因子必须成对出现}


设 \(\beta(x,y)=x^{\mathsf T}B_1y\)。由 \(B_2=B_1A\) 的交替性可得



\[
\beta(Ax,y)=\beta(x,Ay),\qquad \beta(v,Av)=0.
\]



第一式叫 \(A\) 对 \(\beta\) 自伴。证明要找到相互配对的两个同型循环模，并把它们一起正交剥离。



首先，任意 \(j\ge0\)，自伴性给



\[
\beta(v,A^{2j}v)=\beta(A^jv,A^jv)=0,
\]



\[
\beta(v,A^{2j+1}v)=\beta(A^jv,A^{j+1}v)=0.
\]



后一式用 \(\beta(u,Au)=0\)。由双线性，\(\beta(A^iv,A^jv)=0\) 对全部 \(i,j\)；故每个循环模都是全迷向的，即它内部的 \(\beta\) 恒零。它无法单独非退化，必须找一个伙伴。



不同不可约多项式的主分量彼此正交。具体地，若 \(x\) 在 \(f\) 主分量、\(y\) 在 \(g\) 主分量且 \(f\ne g\)，某个 \(f(A)^m\) 杀掉 \(x\)，但在 \(g\) 主分量上可逆。所以写 \(y=f(A)^m y_0\)，就有



\[
\beta(x,y)=\beta(f(A)^mx,y_0)=0.
\]



\Needspace{5\baselineskip}
全空间非退化于是保证每个主分量内部也非退化。



现在固定一个 \(f\) 主分量，选向量 \(v\) 的湮灭多项式恰为这个分量中最大的 \(f^k\)，令 \(X=\mathbb F_2[A]v\)。因为 \(f(A)^{k-1}v\ne0\)，非退化性给同一主分量的 \(w\)，使



\[
\beta(f(A)^{k-1}v,w)\ne0.
\]



令 \(Y=\mathbb F_2[A]w\)。考察配对映射 \(X\to Y^*\)，\(x\mapsto(y\mapsto\beta(x,y))\)。其核 \(K_0\) 是 \(A\) 不变子模：若 \(x\) 与 \(Y\) 全部正交，则 \(\beta(Ax,y)=\beta(x,Ay)=0\)，因为 \(Ay\in Y\)。



为什么这个核一定零？循环模 \(X\cong\mathbb F_2[t]/(f^k)\) 的子模是 \((f^j)/(f^k)\)。任何非零子模都包含最底层



\[
f^{k-1}X,
\]



也称socle（底层简单子模）。它的二元维数是 \(\deg f\)，不一定是1。若 \(K_0\ne0\)，它就包含 \(f(A)^{k-1}v\)，与刚才选的非零配对矛盾。故 \(X\to Y^*\) 单射，\(\dim X\le\dim Y\)。



另一方面 \(k\) 已最大，\(w\) 的湮灭指数不超过 \(k\)，所以 \(\dim Y\le\dim X=k\deg f\)。于是二者等维，配对完美，\(Y\) 也有初等因子 \(f^k\)。两者都全迷向，完美配对又保证 \(X\cap Y=0\)，且 \(X\oplus Y\) 非退化。其正交补因自伴性仍 \(A\) 不变，继续重复，便得到成对循环块的全分解。



此论证不仅说特征多项式是平方；它更强，要求每个具体幂次的初等因子 \(f^k\) 都成对，这正是后面识别 \(p^2\) 核心所需要的。



\Needspace{7\baselineskip}
\section{命题3.4：根空间维数怎样受半维数奇偶限制}


对一个 \(p^k\) 循环模，乘 \(p\) 的核是 \((p^{k-1})/(p^k)\)，维数2。因此一对 \(p^k\) 块占 \(4k\) 维，却只给 \(\ker p(A)\) 贡献4维。其他不可约 \(f\) 块上 \(p(A)\) 可逆，贡献0。由于初等因子成对，立刻得



\[
\rho\equiv0\pmod4.
\]



不存在一次因子，而二元域唯一不可约二次多项式就是 \(p=t^2+t+1\)。所以非 \(p\) 块的次数至少3，一对至少占6维。



\Needspace{5\baselineskip}
把所有对的维数除以2，得到



\[
n=\sum_{\text{成对块}} k\deg f.
\]



若 \(n\) 奇数，至少一项 \(k\deg f\) 为奇数，所以 \(k\) 与 \(\deg f\) 都奇。这样的 \(f\) 不可能是次数2的 \(p\)，又不可能一次，故非 \(p\) 部分至少占6维。剩余维数即使全部由最有效的简单 \(p\) 对组成，也只能给



\[
\rho\le2n-6.
\]



若 \(n\) 偶数，可能全部都是简单 \(p\) 对，这时 \(\rho=2n\)。但若如此，\(C=0\)，\(D\) 是线性且 \(D(0)=0\)，其字符和只能是 \(N\)（零函数）或0（非零线性函数），不可能负。因此\textbf{在要求负字符和时}，\(\rho<2n\)。两者都被4整除，所以



\[
\mathcal G(D)<0\quad\Longrightarrow\quad\rho\le2n-4.
\]



这就是偶数上下界不对称的根源。上端可以令所有点正号；下端若要负号，至少损失4维根空间。它不是某个构造偶然产生的不对称。



\Needspace{7\baselineskip}
\section{从根维数到数值界，每个幂次都核对}


奇数 \(n\ge3\) 时，



\[
|\mathcal G(D)|\le2^{n+(2n-6)/2}=2^{2n-3},
\]



故 \(|\mathcal P|\le2^{3n}2^{2n-3}=2^{5n-3}\)。除以 \(4N=2^{2n+2}\)，英文(3.2)给



\[
M_n-2^{3n-5}\le v_4\le M_n+2^{3n-5}.
\]



偶数 \(n\ge4\) 时，正字符和最大为 \(N=2^{2n}\)，所以 \(\mathcal P\le2^{5n}\)，给上界 \(M_n+2^{3n-2}\)。负字符和只能在 \(\rho\le2n-4\) 出现，故



\[
\mathcal G(D)\ge-2^{2n-2},\quad
\mathcal P\ge-2^{5n-2},\quad
v_4\ge M_n-2^{3n-4}.
\]



若 \(D\) 在根上非零，则 \(\mathcal P=0\)，得到中间基准 \(M_n\)；这不意味着所有中间值都存在。



小维数必须单独处理。\(n=1\) 时 \(V\) 二维，交替形式空间只有1维，只有一个非零交替形式。两分量若都非零便相等，其和为零，不可能三个均非退化。所以这个群类不存在。\(n=2\) 时 \(V\) 四维，成对块只容一对简单 \(p\) 块，故 \(C=0\)。\(D\) 为线性函数：恒零给 \(M_2+2^{4}=69\)，非零给 \(M_2=53\)，没有负字符和。因此不能把一般偶数下界公式硬套到 \(n=2\)。



\Needspace{7\baselineskip}
\section{命题3.5的全部充要等号条件}


\Needspace{5\baselineskip}
界来自多个不等式。等号必须同时让每一步取等，不能只要求某个矩阵有大根空间。



\textbf{奇数 \(n\ge3\) 的任一端点：} 必须 \(\rho=2n-6\)。用“亏损” \(2n-\rho\) 检查块结构：一对 \(p^k\) 块的亏损是 \(4k-4=4(k-1)\)，其他对的亏损是整块维数 \(2k\deg f\ge6\)。奇数性已经迫使至少一个非 \(p\) 对，亏损总共只有6，因此恰好一个简单不可约三次对，其他全为简单 \(p\) 对。再要求 \(D|_{\ker C}=0\)，它才能下降到6维非退化商；该商字符和应为下端 \(-8\) 或上端 \(+8\)。乘回根空间大小 \(2^{2n-6}\) 后，正好给所需全局字符和。反之，这三项一起成立，逐式代回就达到对应端点。



\textbf{偶数 \(n\ge4\) 的下端：} 必须 \(\rho=2n-4\)，即亏损4。不可能有非 \(p\) 对，因为那至少亏6。\(\sum4(k-1)=4\) 强迫恰好一个 \(p^2\) 对，其余全为简单 \(p\) 对。再要求 \(D\) 在根上为零，以及4维非退化商的字符和为 \(-4\)。这些条件既必要也充分。这是后来八维核心的结构来源：一对 \(p^2\) 模各4维，共8维；并不是先看到某个八维例子，再假定所有等号群都含它。



\textbf{偶数 \(n\ge2\) 的上端：} 必须 \(\mathcal G(D)=N\)。这是 \(N\) 个 \(\pm1\) 相加，达到 \(N\) 当且仅当每项为 \(+1\)，即 \(D\equiv0\)。这也迫 \(C=0\) 和 \(p(A)=0\)；但反过来只知道 \(p(A)=0\) 还需剩余线性函数 \(D\) 恒零，才真到上端。



\textbf{\(n=2\) 的下端：} \(C=0\) 已自动成立；下端恰是 \(D\) 为非零线性函数。不能把它与高维“根上消失且负商符号”混为一谈。



这些充要条件完成数值界的逻辑证明，但仍有两件事未完成：端点是否真的存在，以及偶数下端的完整平方映射是否只有一个严格等价类。下一章处理第一件；八维细化与高维相位章节处理第二件，不能用尚待证明的分类反过来证明本章界成立。



\chapter{每个端点都真实存在：四维、六维与八维构造}


本章对应英文第4节，基线第7—9页：引理4.1、4.2、4.3及定理1.1的完成证明。上一章给的是“任何群不能越过的界”。本章必须构造函数并核对：三个标量极化非退化、三个原始Gauss和全正、第三矩有精确数值。命题2.1才能据此造出满足全部群假设的群。构造不依赖后面的唯一性分类，因而避免循环论证。



\Needspace{7\baselineskip}
\section{一个会反复使用的求和方法}


令 \(L=\mathbb F_{2^m}\)，\(\operatorname{Tr}_L\) 为到 \(\mathbb F_2\) 的绝对迹，且 \(a\ne0\)。考虑



\[
\sum_{x,y\in L}(-1)^{\operatorname{Tr}_L(axy+cx+dy+ux+vy)}.
\]



对固定 \(x\) 先对 \(y\) 求和，\(y\) 的系数是 \(ax+d+v\)。迹配对非退化，所以该内和仅在 \(ax+d+v=0\) 时非零，非零值是 \(2^m\)。唯一满足的 \(x\) 为 \((v+d)/a\)，故总和等于



\[
2^m(-1)^{\operatorname{Tr}_L((u+c)(v+d)/a)}.
\tag{英文(4.2)，在本章所用的维数下}
\]



这既算出了原始Gauss符号（令 \(u=v=0\)），又算出全部实际Walsh对偶（保留 \(u,v\)）。它是一般Maiorana–McFarland公式[CMP18，式(4)，印刷页313]的直接实例；任意辅助函数 \(g\) 的一般式会在后续相位计算及附录中用到。当前不把这条经典求和方法说成新方法。



\Needspace{7\baselineskip}
\section{引理4.1：四维正端块 \texorpdfstring{\(H_+\)}{H+}}


取 \(K=\mathbb F_4\)，\(\alpha^2+\alpha+1=0\)，\(\operatorname{Tr}(w)=w+w^2\)。有



\[
\operatorname{Tr}(0)=\operatorname{Tr}(1)=0,\qquad
\operatorname{Tr}(\alpha)=\operatorname{Tr}(\alpha^2)=1.
\]



\Needspace{5\baselineskip}
英文(4.1)用



\[
\iota(w)=(\operatorname{Tr}(w),\operatorname{Tr}(\alpha w))
\]



把 \(W=K\) 识别为 \(\mathbb F_2^2\)。非零分量可统一写作 \(q_z=\operatorname{Tr}(zQ)\)，其中 \(z\in K^\times\)；\(z=1,\alpha,\alpha^2\) 恰对应三个非零泛函。这是一组共同的目标坐标，后面拼接各块始终保留它。



对 \(H_+(u,v)=uv\)，极化为 \(uv'+u'v\)，标量极化是 \(\operatorname{Tr}(z(uv'+u'v))\)。若 \(u\ne0\)，可选 \(v'\) 使 \(\operatorname{Tr}(zuv')\ne0\)；若 \(u=0,v\ne0\) 则交换角色，所以每个 \(z\ne0\) 都非退化。



频率变量使用同一迹配对 \(\operatorname{Tr}(\rho u+\sigma v)\)。上一节求和给



\[
\sum_{u,v\in K}(-1)^{\operatorname{Tr}(zuv+\rho u+\sigma v)}
=4(-1)^{\operatorname{Tr}(\rho\sigma/z)}.
\]



零频都为 \(+4\)，所以该块达到最大 \(s\)。三个对偶之和为



\[
D(\rho,\sigma)=\operatorname{Tr}\left(\rho\sigma\sum_{z\in K^\times}z^{-1}\right)=0,
\]



因为 \(K^\times\) 的逆仍是 \(1,\alpha,\alpha^2\)，和为0。因此 \(\mathcal G(D)=16\)，\(h=4\)，



\[
\mathcal P=4^3\cdot16=2^{10}.
\]



这不是依靠“正端”名称假定符号，而是完整计算。群阶为 \(4\cdot16=64\)，\(s=16+3\cdot4=28\)。



零点也很直观：\(uv=0\) 当且仅当 \(u=0\) 或 \(v=0\)，所以两条 \(K\) 坐标轴各4点、相交于零，共 \(z=7\) 点。它们也是仅有的两个二维二元全零子空间。若一个全零子空间含两条轴各自的非零向量，它们的和两坐标非零，乘积非零，矛盾。于是 \(L=2\)，第2章公式给



\[
v_4=16\cdot2+6\cdot7-5=69.
\]



\Needspace{7\baselineskip}
\section{同极化但不同群：\texorpdfstring{\(uv+u\)}{uv+u} 给 \texorpdfstring{\(n=2\)}{n=2} 下端}


考虑 \(Q'(u,v)=uv+u=u(v+1)\)。添加一次项不改变极化，所以三个标量极化仍非退化。其 \(z\) 分量是在上一求和式中把 \(\rho\) 替换为 \(\rho+z\)，故零频仍全正；实际对偶为



\[
q_z^*(\rho,\sigma)=\operatorname{Tr}((\rho+z)\sigma/z).
\]



三个相加，含 \(\rho\sigma\) 的部分仍消掉，而每项留下 \(\operatorname{Tr}(\sigma)\)，三次相加在特征2里还是它。因此



\[
D(\rho,\sigma)=\operatorname{Tr}(\sigma),
\]



非零线性、极化为零、整数字符和为零。这正好实现命题3.5的 \(n=2\) 下端条件。



直接零点计数也可独立核对。零集由 \(u=0\) 的4点和 \(v=1,u\ne0\) 的3点组成，共7点。唯一全零二维二元子空间是 \(u=0\)：如果含两个不同的第二类点 \((u,1),(u',1)\)，其和 \((u+u',0)\) 不为零点；如果含一个第二类点与非零的 \((0,v)\)，其和的 \(Q'\) 值为 \(uv\ne0\)。所以 \(L=1\)，



\[
v_4=16+42-5=53.
\]



两个映射拥有相同的极化、相同的 \(s=28\)，但 \(v_4\) 不同，故对应群不同构。这是全文最小、最具体的例子，说明不能在平方映射分类时任意抹去线性项。此处用零点计数与字符和两种方法算同一个结果，是对转换接口的核对，不是另一个原创主张。



\Needspace{7\baselineskip}
\section{引理4.2：为什么奇数分支需要一个六维块}


奇数半维数端点的极化必须含一对简单不可约三次块。因此自然用 \(L=\mathbb F_8\) 的乘法。写



\[
L=\mathbb F_2[t]/(t^3+t+1),\qquad t^3=t+1,
\]



且 \(\operatorname{Tr}_8(w)=w+w^2+w^4\)。在 \(L^2\)（二元维数6）上取两分量



\[
q_1(x,y)=\operatorname{Tr}_8(xy),\qquad
q_2(x,y)=\operatorname{Tr}_8(txy+cx+dy).
\]



第三分量是它们的函数和，二次乘积系数为 \(1+t\)。\(1,t,1+t\) 全非零，所以三个标量极化都非退化。此处输出仍是两维二元空间；没有把 \(L\) 本身误当中心，也没有把中心扩到8个元素。



\Needspace{5\baselineskip}
先列会用到的幂，便于逐项核查：



\[
\begin{array}{c|ccccc}
j&3&4&5&6&7\\\hline
t^j&t+1&t^2+t&t^2+t+1&t^2+1&1
\end{array}
\]



所以 \(t^{-1}=t^6\)，\((1+t)^{-1}=t^4\)，并且



\[
k=t^{-1}+(1+t)^{-1}=t+1,\qquad 1+k=t.
\]



迹值为 \(\operatorname{Tr}_8(1)=1\)、\(\operatorname{Tr}_8(t)=\operatorname{Tr}_8(t^2)=0\)。



当 \(c=d=0\)，三个原始和由求和式都是 \(+8\)。实际对偶相加得



\[
D(u,v)=\operatorname{Tr}_8\bigl((1+t^{-1}+(1+t)^{-1})uv\bigr)=\operatorname{Tr}_8(tuv).
\]



这是非退化六维二次函数，字符和为 \(+8\)。故 \(\mathcal P=8^3\cdot8=2^{12}\)，实现奇数分支最小维数 \(n=3\) 的上端。



\Needspace{7\baselineskip}
\section{六维负端：逐项保留原始相位}


\Needspace{5\baselineskip}
现在取



\[
c=1,\qquad d=t^2+t+1=t^5.
\]



\Needspace{5\baselineskip}
不能因为极化未变就断言原始Gauss和仍正；需检查两项可能的原点相位：



\[
\operatorname{Tr}_8(cd/t),\qquad\operatorname{Tr}_8(cd/(1+t)).
\]



由幂表，\(d/t=t^4=t^2+t\)，\(d/(1+t)=t^2\)，迹均为0。因此三个原始和确实都是 \(+8\)。按同一频率配对，三个实际对偶完整相加为



\[
D(u,v)=\operatorname{Tr}_8(tuv+kdu+kcv+kcd).
\]



这里的常数 \(kcd\) 是对偶公式实际给出的项；即使最后其迹为零，也应在代入前保留，避免把块级原点相位混淆。应用求和式，



\[
\mathcal G(D)=8(-1)^{\operatorname{Tr}_8(cd(k+k^2/t))}.
\]



因为 \(k=t+1\)，特征2下 \(k^2=t^2+1\)，所以



\[
k+k^2/t=k/t=t^2.
\]



而 \((k/t)d=t^2t^5=t^7=1\)，迹是1，故



\[
\mathcal G(D)=-8,\qquad \mathcal P=8^3(-8)=-2^{12}.
\]



\(D\) 的极化由系数 \(t\) 给出，依然非退化。这个例子把“原来的三个和全正”与“它们实际对偶相加后的和为负”清楚分开：二者并不矛盾。数值上 \(M_3=447\)，偏移 \(2^{3\cdot3-5}=16\)，得到上下端431与463。



\Needspace{7\baselineskip}
\section{引理4.3：八维负端为什么使用 \texorpdfstring{\(p^2\)}{p2} 局部环}


偶数下端等号要求一对 \(p^2\) 块，\(p(t)=t^2+t+1\)。取局部环



\[
R=\mathbb F_2[t]/(p(t)^2),
\]



它的二元维数是4。令 \(U\) 为乘 \(t\) 的线性变换，\(E=p(U)=U^2+U+I\)。在这个环中 \(p^2=0\)，故 \(E^2=0\)；但 \(E\ne0\)。乘 \(p\) 的像和核都恰好是理想 \((p)\)，维数2，所以



\[
\operatorname{rank}E=2,\qquad \operatorname{im}E=\ker E.
\]



这是非半单块：存在非零幂零部分 \(E\)。与两个简单 \(p\) 块的直和相比，后者满足 \(p(U)=0\)，这里却不满足；特征多项式相同不代表等价。



选 \(0\ne c\in\operatorname{im}E\)，再选列向量 \(z\) 使 \(z^{\mathsf T}c=1\)，置 \(d=E^{\mathsf T}z\)。这些选择总存在：非零向量总能被某个线性泛函取到1。\(c\in\operatorname{im}E\) 又给 \(Ec=0\)。



\Needspace{5\baselineskip}
在两个四维二元空间上定义



\[
q_1(x,y)=x^{\mathsf T}y,\qquad
q_2(x,y)=x^{\mathsf T}Uy+c^{\mathsf T}x+d^{\mathsf T}y.
\]



相应三极化的乘积系数是 \(I,U,I+U\)。由于 \(t\) 与 \(t+1\) 都和 \(p^2\) 互素，它们在 \(R\) 中是单位，故 \(U,I+U\) 可逆。这证明三个标量极化非退化。



\textbf{教学补充例：原构造的一组具体坐标。} 在多项式基 \(1,t,t^2,t^3\) 下有



\[
U=\begin{pmatrix}0&0&0&1\\1&0&0&0\\0&1&0&1\\0&0&1&0\end{pmatrix},
\]



因为 \(t^4=t^2+1\)。可取 \(c=(1,1,1,0)^{\mathsf T}=E(1,0,0,0)^{\mathsf T}\)、\(z=(1,0,0,0)^{\mathsf T}\)、\(d=(1,0,1,1)^{\mathsf T}\)。这是原构造的一个坐标实现，仅为便于直接核查，并不声称带来了新的群类型。



\Needspace{7\baselineskip}
\section{八维块的原始正号及完整对偶和}


一般的二元矩阵求和与有限域求和一样：若 \(M\) 为可逆四阶矩阵，则



\[
\sum_{x,y\in\mathbb F_2^4}
(-1)^{x^{\mathsf T}My+c^{\mathsf T}x+d^{\mathsf T}y+a^{\mathsf T}x+b^{\mathsf T}y}
=16(-1)^{(b+d)^{\mathsf T}M^{-1}(a+c)}.
\]



令频率为零，两项可能的原始相位是 \(d^{\mathsf T}U^{-1}c\) 与 \(d^{\mathsf T}(I+U)^{-1}c\)。因 \(d^{\mathsf T}=z^{\mathsf T}E\)，\(E\) 与 \(U\) 及其可逆多项式都交换，且 \(Ec=0\)，这两项都为0。例如



\[
d^{\mathsf T}U^{-1}c=z^{\mathsf T}EU^{-1}c=z^{\mathsf T}U^{-1}Ec=0.
\]



所以三个原始和都是 \(+16\)。至此才可使用最大 \(s\) 分支下无原点相位的全局对偶记号。



令 \(K_0=U^{-1}+(I+U)^{-1}\)。因为



\[
U(I+U)=U+U^2=I+E,\qquad (I+E)^{-1}=I+E,
\]



有 \(K_0=(U(I+U))^{-1}=I+E\)，所以



\[
I+K_0=E,\quad K_0c=c,\quad EK_0=E.
\]



频率 \((a,b)\) 下，三个实际对偶相加，先不删线性和常数项，得到



\[
\begin{aligned}
D(a,b)
&=b^{\mathsf T}a+(b+d)^{\mathsf T}U^{-1}(a+c)
 +(b+d)^{\mathsf T}(I+U)^{-1}(a+c)\\
&=b^{\mathsf T}(I+K_0)a+b^{\mathsf T}K_0c
  +d^{\mathsf T}K_0a+d^{\mathsf T}K_0c\\
&=b^{\mathsf T}(Ea+c)+z^{\mathsf T}Ea.
\end{aligned}
\]



最后一行使用 \(d^{\mathsf T}K_0=z^{\mathsf T}EK_0=z^{\mathsf T}E\)，且常数 \(d^{\mathsf T}K_0c=z^{\mathsf T}Ec=0\)。这说明原证明中的简式包含了全部线性与常数项的真实消去过程，不是任意减去对偶常数。



现在求和时先对 \(b\)。若 \(Ea+c\ne0\)，内和为0；若 \(Ea=c\)，内和为16。由于 \(c\in\operatorname{im}E\)，这个方程有 \(|\ker E|=4\) 个解。在这些解上



\[
z^{\mathsf T}Ea=z^{\mathsf T}c=1.
\]



\Needspace{5\baselineskip}
所以每个解给负号，



\[
\mathcal G(D)=-16\cdot4=-64,\qquad
\mathcal P=16^3(-64)=-2^{18}.
\]



还可直接核对等号的根条件：\(D\) 的二次部分是 \(b^{\mathsf T}Ea\)，极化矩阵为 \(\begin{pmatrix}0&E^{\mathsf T}\\E&0\end{pmatrix}\)，秩4、根维数4。在根上有 \(Ea=0,E^{\mathsf T}b=0\)；又 \(c\in\operatorname{im}E\)，所以 \(b^{\mathsf T}c=0\)，即 \(D\) 在根上为零。商四维的字符和是 \(-64/2^4=-4\)，恰符合命题3.5(b)。



因此这个8维映射实现 \(n=4\) 的下端。此刻无需知道它与第1章写出的 \(Q_*\) 严格等价；那是后面完整分类将证明的内容。把可达性与分类分开，是证明组织的重要逻辑边界。



\Needspace{7\baselineskip}
\section{用共同目标正交和得到全部允许维数}


若 \(Q=Q_0\perp Q_1\)，对同一个目标泛函 \(\lambda\) 有 \(q_\lambda(x_0,x_1)=q_{\lambda,0}(x_0)+q_{\lambda,1}(x_1)\)。频率也按直和分解，有限和直接拆开：



\[
\mathcal W_\lambda(\xi_0,\xi_1)
=\mathcal W_{\lambda,0}(\xi_0)\mathcal W_{\lambda,1}(\xi_1).
\]



\Needspace{5\baselineskip}
所以原始Gauss和相乘。对三个分量乘积再对全部频率求和，仍可分离，得到



\[
\mathcal P(Q_0\perp Q_1)=\mathcal P(Q_0)\mathcal P(Q_1).
\]



共同目标很重要：这不是允许各块独立重新标号后再不加说明地相乘。当前每块都已经用同一有序二维目标构造，故可以直接拼接。



\textbf{偶数上端。} \(n/2\) 个 \(H_+\) 块共有 \(4(n/2)=2n\) 维，原始和乘成 \(4^{n/2}=2^n=h\)，而



\[
\mathcal P=(2^{10})^{n/2}=2^{5n}.
\]



这也是[CMP18, Theorem 2(i)，印刷页314]的标准迹乘积构造：取输入域 \(\mathbb F_{2^n}\)、子域 \(K=\mathbb F_4\)、\(\pi(y)=y\)、\(g=0\)，映射 \(\operatorname{Tr}_{\mathbb F_{2^n}/K}(xy)\) 在一对迹对偶的 \(K\) 基中就是 \(H_+^{\perp n/2}\)。\(2\mid n\) 是子域存在所需条件，不能在奇数 \(n\) 无条件使用。



\textbf{奇数两个端点。} 取六维正块或负块，再加 \((n-3)/2\) 个 \(H_+\)。维数为 \(6+4(n-3)/2=2n\)，原始和乘成 \(8\cdot4^{(n-3)/2}=h\)，且



\[
\mathcal P=\pm2^{12}\,2^{10(n-3)/2}=\pm2^{5n-3}.
\]



\textbf{偶数下端。} 取八维负块，再加 \((n-4)/2\) 个 \(H_+\)，维数为 \(2n\)，原始和是 \(16\cdot4^{(n-4)/2}=h\)，而



\[
\mathcal P=-2^{18}\,2^{10(n-4)/2}=-2^{5n-2}.
\]



\textbf{\(n=2\)。} 已由 \(uv\) 与 \(uv+u\) 分别给69和53。



每个构造的标量极化都为非退化块的正交和，因而仍非退化；命题2.1给实际semi-extraspecial群，且 \(Z=G'=\Phi\cong C_2^2\)。原始和全正保证最大 \(s\)，\(\mathcal P\) 的数值配合英文(3.2)保证恰到所述端点。这就同时证明群类在每个 \(n\ge2\) 非空，以及定理1.1全部界是最优且可达。



接下来要解决的不是再扩大构造，而是反向穷尽：任何偶数下端等号映射，是否都能化成已指定的八维 \(Q_*\) 与正余块？上一章已经把极化限制到一个 \(p^2\) 对及若干简单 \(p\) 对；后续必须保留每个线性细化，才能从这个极化结构走到真正的群同构结论。



\chapter{八维核心：从全部细化到一个严格同构类型}


本章逐步解释原稿§5：引理5.1（基线第9—10页）、引理5.2（第10页）、引理5.3与命题5.4、5.5（第11—12页）。这里“八维”始终指商空间 \(V=G/Z(G)\) 在 \(\mathbb F_2\) 上的维数；相应群阶为 \(4\cdot2^8=1024\)。本章只分类已经满足最大自逆元素数和下界等号的八维对象。不能把固定某种极化后的计算，误当作所有八维平方映射的分类。



上一章已经构造了一个下界对象，证明下界可达。本章的问题更强：还有没有其他群同构类型？答案需要两个独立步骤：先覆盖全部可能的平方映射，再用允许的换基把全部筛选结果化到同一代表。只比较交换子映射，或只证明一个构造满足等号，都不够。



\Needspace{7\baselineskip}
\section{等号先告诉我们什么}


沿用第三章的记号：\(B_1,B_2\) 是两个标量极化矩阵，\(A=B_1^{-1}B_2\)，\(p(t)=t^2+t+1\)。原稿命题3.5(b)已经给出充要条件的一部分：偶数半维数下端的极化必须恰有一对 \(p^2\) 循环块，其余为简单 \(p\) 块对。在八维情形没有余下空间，因此恰是两个相同的 \(p^2\) 循环模。



“一对 \(p^2\) 块”与“两对简单 \(p\) 块”不同。前者每个循环块的最小多项式是 \(p^2\)；后者的最小多项式只是 \(p\)。两种情形的特征多项式都可能是 \(p^4\)，但只有前者非半单。不能仅看特征多项式便将它们合并。



由引理3.3，我们能选出互相配对的两个循环子空间；每个单独对 \(B_1\) 都是全迷向的，但两者之间的配对非退化。把算子 \(A\) 的作用写成“乘 \(t\)”，这两个空间都成为



\[
R=\mathbb F_2[t]/(p^2)
\]



的一个自由秩一模，合起来为 \(R^2\)。\(R\) 有 \(2^4=16\) 个元素，\(R^2\) 才是八维二元空间。这里的环不是一个有16个元素的域；它含有非零幂零元。



更高维时，同样先取这个 \(R^2\) 作为 \(V_0\)，再取 \(B_1\) 正交补 \(V_1\)。因为 \(A\) 对 \(B_1\) 自伴，若 \(x\in V_0\)、\(y\in V_1\)，则



\[
B_1(x,Ay)=B_1(Ax,y)=0.
\]



所以 \(V_1\) 对 \(A\) 不变，并同时与 \(V_0\) 关于 \(B_1,B_2\) 正交。于是这是同一个向量值极化 \(B\) 的共同正交分解。这个步骤只用极化，没有推出任何限制块的原始Gauss和为正；符号要到下一章才处理。



\Needspace{7\baselineskip}
\section{为什么局部环可写成四元域加一个幂零方向}


在 \(R\) 中定义



\[
\varepsilon=p(t)=t^2+t+1,\qquad \alpha=t+\varepsilon=t^2+1.
\]



商去 \(p^2\) 正是说 \(\varepsilon^2=0\)。又有



\[
\alpha^2+\alpha+1
=(t+\varepsilon)^2+t+\varepsilon+1
=p(t)+\varepsilon=0.
\]



因此 \(\alpha\) 生成一个子域 \(K\cong\mathbb F_4\)，其中 \(\alpha^2=\alpha+1\)、\(\alpha^3=1\)。每个环元唯一写成 \(a+\varepsilon b\)，\(a,b\in K\)，故



\[
R=K[\varepsilon]/(\varepsilon^2).
\]



这是将一个四维二元计算压缩成两个四元域坐标的办法，并未把原平方映射变成 \(K\) 上齐次二次型。\(\varepsilon\) 记录了不能被域坐标消除的非半单部分。环元 \(a+\varepsilon b\) 可逆当且仅当 \(a\ne0\)：若 \(a\ne0\)，其逆为 \(a^{-1}+\varepsilon b a^{-2}\)；若 \(a=0\)，其平方为零。



令 \(\lambda:R\to\mathbb F_2\) 提取多项式代表中的 \(t^3\) 系数。配对 \((x,y)\mapsto\lambda(xy)\) 非退化：非零代表 \(x\) 的次数若为 \(j\le3\)，乘以 \(t^{3-j}\) 后恰有最高项 \(t^3\)，还没有超过三次而发生模约化，故 \(\lambda(xt^{3-j})=1\)。



为什么可以把实际配对规范到这一配对？因为循环模之间的平衡配对由一个线性泛函 \(\lambda'\) 决定，而 \(e\mapsto[z\mapsto\lambda(ez)]\) 是两个四维二元空间间的线性同构。因此存在唯一 \(e\in R\) 使 \(\lambda'(z)=\lambda(ez)\)。实际配对非退化迫 \(e\) 可逆；否则乘 \(e\) 有非零核，会成为配对的根。将一个循环生成元乘以 \(e^{-1}\) 即完成规范化，没有变动目标坐标。



\Needspace{5\baselineskip}
于是原稿式(5.3)的两个标量极化成为



\[
B((x,y),(x',y'))=
\bigl(\lambda(xy'+x'y),\lambda(t(xy'+x'y))\bigr).
\]



这里第一项是 \(B_1\)，第二项是 \(B_2\)；同一个 \(t\) 在两个循环模上起作用。对应简单 \(p\) 块时，\(t\) 就变成乘以 \(\alpha\)，可用同一目标顺序的迹配对规范化。



\Needspace{7\baselineskip}
\section{从环配对得到全部平方映射}


记 \(\operatorname{Tr}=\operatorname{Tr}_{K/\mathbb F_2}\)，即 \(\operatorname{Tr}(w)=w+w^2\)。四个值为



\[
\operatorname{Tr}(0)=\operatorname{Tr}(1)=0,
\qquad \operatorname{Tr}(\alpha)=\operatorname{Tr}(\alpha^2)=1.
\]



在基 \(1,\alpha,\varepsilon,\alpha\varepsilon\) 上直接计算，得



\[
\lambda(a+\varepsilon b)=\operatorname{Tr}(b),\qquad
\lambda(t(a+\varepsilon b))=\operatorname{Tr}(a+\alpha b).
\]



用 \(x=a+\varepsilon b\)、\(y=c+\varepsilon d\) 写环坐标，则 \(xy=ac+\varepsilon(ad+bc)\)。基准平方映射的标量值为



\[
\bigl(\operatorname{Tr}(ad+bc),\operatorname{Tr}(ac+\alpha(ad+bc))\bigr).
\]



原稿固定目标识别 \(\iota(w)=(\operatorname{Tr}(w),\operatorname{Tr}(\alpha w))\)。若将上述有序对写回 \(K\)，其值正是



\[
Q_{\mathrm{base}}=ad+bc+\operatorname{Tr}(ac).
\]



可核对第二分量：\(\operatorname{Tr}(\alpha\operatorname{Tr}(ac))=\operatorname{Tr}(ac)\)，因为 \(\operatorname{Tr}(ac)\in\mathbb F_2\) 且 \(\operatorname{Tr}(\alpha)=1\)。



现在必须补齐所有与基准映射有同一极化的映射。若 \(Q,Q_0\) 的极化相同，\(L=Q+Q_0\) 满足 \(L(x+y)=L(x)+L(y)\)；在二元空间上这恰是线性。因此全部细化是 \(Q_0+L\)，不能只选最方便的一种。



每个二元线性映射 \(K\to K\) 都唯一写成 \(a\mapsto fa+ga^2\)，\(f,g\in K\)。这是因为 \(a\mapsto a^2\) 在 \(\mathbb F_2\) 上线性；若 \(fa+ga^2\) 恒零，代 \(a=1\) 得 \(f=g\)，再代 \(a=\alpha\) 得 \(f(\alpha+\alpha^2)=f=0\)。这给一个四维二元空间到另一个四维二元空间的单射，所以也是满射。



独立作用于 \(a,b,c,d\) 四个坐标，便得到原稿式(5.1)的完整列表：



\[
\begin{aligned}
Q^{(0)}={}&ad+bc+\operatorname{Tr}(ac)
+fa+ga^2+\ell c+mc^2\\
&+\eta b+kb^2+\theta d+jd^2,
\end{aligned}
\]



八个系数都在 \(K\)，共有 \(4^8=2^{16}\) 个带坐标映射。这里 \(ga^2\) 虽然写成平方，却也是二元线性项；\(fa\) 虽然是普通一次项，却不能被任意删掉。所有这些映射共有极化，未必有相同的零点数、四元子群数或群同构类型。



\Needspace{7\baselineskip}
\section{全部核心相位：先算Walsh和，不先假定正号}


引理5.2评价所有上述候选。三个非零标量分量用 \(z\in K^\times\) 标记，取 \(q_z=\operatorname{Tr}(zQ^{(0)})\)。频率为



\[
\operatorname{Tr}(\rho a+\sigma b+\tau c+\upsilon d).
\]



这个非退化迹配对表示全部二元线性泛函，没有遗漏频率。令 \(\kappa_z=\operatorname{Tr}(z)\)。先用



\[
\operatorname{Tr}(zb^2k)=\operatorname{Tr}(z^2k^2b)
\]



把平方项移入 \(b\) 的线性系数。恒等式来自 \(\operatorname{Tr}(w)=\operatorname{Tr}(w^2)\) 及 \(b^4=b\)。因此在固定 \(a,c\) 后，\(b\) 的系数为



\[
zc+z\eta+z^2k^2+\sigma.
\]



对 \(b\in K\) 求和：这个系数非零时和为零，系数为零时和为4。同理，\(d\) 的系数为 \(za+z\theta+z^2j^2+\upsilon\)。故全部求和只留下



\[
a=a_z=\upsilon/z+\theta+zj^2,
\qquad c=c_z=\sigma/z+\eta+zk^2,\tag{\text{原稿 5.4}}
\]



并产生 \(4\cdot4=16\) 倍率。剩余指数给



\[
\mathcal W_z=16(-1)^{E_z},
\]



\[
E_z=\operatorname{Tr}\left(
\kappa_z a_zc_z+(\rho+zf+z^2g^2)a_z
+(\tau+z\ell+z^2m^2)c_z\right).\tag{\text{原稿 5.5}}
\]



尤其 \(E_z(0)\) 是原始Gauss和的符号：为0才是正16，为1则是负16。不能先把它减掉，再声称三个块都为正。此直接求和也是CMP2018式(4)的坐标代入；附录讲解会说明三个分量使用的共同 \(\mathbb F_{16}\) 坐标。



将三个 \(E_z\) 在 \(\mathbb F_2\) 中相加，可以把许多项消去。为让式(5.6)可独立核查，列出实际需要的有限和：\(K^\times=\{1,\alpha,\alpha^2\}\)，所以



\[
\sum z=\sum z^2=\sum z^{-1}=0,\quad
\sum z^3=1,\quad \sum\kappa_z=0,
\]



\[
\sum\kappa_z/z=\sum\kappa_z/z^2
=\sum\kappa_z z=\sum\kappa_z z^2=1.
\]



例如 \(\kappa_1=0\)、\(\kappa_\alpha=\kappa_{\alpha^2}=1\)，故 \(\sum\kappa_z/z=\alpha^2+\alpha=1\)。这些和在特征2的域内进行，三个1之和是1。



展开 \(\kappa_z a_zc_z\) 后，求和保留下来的是



\[
\sigma\upsilon+\eta\upsilon+\theta\sigma
+\theta k^2+\eta j^2+j^2k^2.
\]



展开 \((\rho+zf+z^2g^2)a_z\) 后保留下来的是 \(\theta\rho+f\upsilon+g^2j^2\)；另一项同理给 \(\eta\tau+\ell\sigma+m^2k^2\)。再用迹不随平方改变，得到



\[
\begin{aligned}
D_0:=\sum_{z\ne0}E_z
={}&\operatorname{Tr}\bigl(\sigma\upsilon+(f+\eta)\upsilon
+(\ell+\theta)\sigma+\theta\rho+\eta\tau\bigr)+\gamma,\\
\gamma={}&\operatorname{Tr}(\theta k^2+\eta j^2+jk+gj+mk)
=\sum_{z\ne0}E_z(0).
\end{aligned}\tag{\text{原稿 5.6—5.7}}
\]



这里 \(D_0\) 是相位函数，不是整数字符和。作为区分的教学例，八维空间上的常数1函数的字符和是 \(-256\)，而不是1或0；这不是说本段固定极化的 \(D_0\) 能成为常数函数，其极化秩始终为4。下一章中 \(D_0\) 还可能含不能先删的常数 \(\gamma\)。



\Needspace{7\baselineskip}
\section{各向异性为什么能阻止某些原点符号}


引理5.3先作一个条件性推论。设 \(\eta=\theta=0\)，并且三个原始Gauss和具有同一个符号 \((-1)^\delta\)。这里尚允许 \(\delta=1\)。在零频处 \(a_z=zj^2,c_z=zk^2\)。置



\[
e=jk,\qquad u=fj^2+\ell k^2,\qquad v=gj+mk.
\]



\Needspace{5\baselineskip}
从式(5.5)得到



\[
E_z(0)=\operatorname{Tr}(z)\operatorname{Tr}(ze)
+\operatorname{Tr}(z^2u)+\operatorname{Tr}(v).
\]



其中 \(z^3=1\) 把 \(g,m\) 项变成常数；\(\operatorname{Tr}(z^2e^2)=\operatorname{Tr}(ze)\) 处理了第一项。为核查原稿的再化简，可在 \(K\) 内展开



\[
\begin{aligned}
\operatorname{Tr}(z)\operatorname{Tr}(ze)
&=(z+z^2)(ze+z^2e^2)\\
&=z^2e+ze^2+e+e^2\\
&=\operatorname{Tr}(z^2e)+\operatorname{Tr}(e).
\end{aligned}
\]



\Needspace{5\baselineskip}
故



\[
E_z(0)=\operatorname{Tr}(z^2(u+e))+
\operatorname{Tr}(v+e).
\]



若 \(u+e\ne0\)，第一项是 \(z\) 的一个非零二元线性泛函，在三个非零域元上不可能为常数：非零泛函的核有两个元素，其中恰一个非零，另外两个非零域元取值1。因此三个相位全相同恰好等价于



\[
fj^2+\ell k^2=jk,
\qquad \operatorname{Tr}(gj+mk+jk)=\delta.\tag{\text{原稿 5.8}}
\]



现在加上 \(\operatorname{Tr}(f\ell)=1\)。它推出 \(f,\ell\ne0\)，且 \(K\) 上二元齐次型



\[
fX^2+XY+\ell Y^2
\]



没有非零零点，也就是各向异性。证明不必背诵Arf分类：若 \(Y=0\)，方程迫 \(X=0\)；若 \(Y\ne0\)，令 \(w=fX/Y\)，将方程乘以 \(f/Y^2\) 后得到 \(w^2+w+f\ell=0\)。但 \(\operatorname{Tr}(w^2+w)=0\)，与 \(\operatorname{Tr}(f\ell)=1\) 矛盾。



把 \((X,Y)=(j,k)\) 代入，式(5.8)便迫 \(j=k=0\)；第二式随即迫 \(\delta=0\)。这是一条条件清楚的限制：各向异性是在下端负字符和出现后才得到的，不能仅由整体最大自逆元素数便引用它。



\Needspace{7\baselineskip}
\section{为什么恰好96个，而且没有漏掉其他细化}


命题5.4把上面几步按正确顺序联立。八维下界等号要求 \(s=304\)、\(v_4=4371\)。其必要性如下。



首先，命题3.5已经保证所有对象的极化都是本章这一对 \(p^2\) 块，所以式(5.1)确实覆盖全部对象。其次，最大 \(s\) 意味着三个 \(E_z(0)\) 都是0，从而 \(\gamma=0\)。式(5.6)的二次部分只有 \(\operatorname{Tr}(\sigma\upsilon)\)，故其极化根恰为 \(\rho,\tau\) 两个 \(K\) 方向，也就是四维二元根。\(D_0\) 在根上消失要求



\[
\operatorname{Tr}(\theta\rho+\eta\tau)=0
\quad\text{对所有 }\rho,\tau\in K.
\]



迹配对非退化，于是 \(\theta=\eta=0\)。这里不能只检查 \(\rho=\tau=1\)，那会漏掉迹为零但非零的系数；必须对全部域元成立。



此时对 \(D_0\) 求整数字符和，\(\rho,\tau\) 自由给出16倍；再先对 \(\sigma\) 求和：



\[
\begin{aligned}
\sum_{\rho,\sigma,\tau,\upsilon}(-1)^{D_0}
&=16\sum_{\sigma,\upsilon}
 (-1)^{\operatorname{Tr}(\sigma\upsilon+f\upsilon+\ell\sigma)}\\
&=16\cdot4\,(-1)^{\operatorname{Tr}(f\ell)}\\
&=64(-1)^{\operatorname{Tr}(f\ell)}.
\end{aligned}\tag{\text{原稿 5.10}}
\]



第二行的原因是：对 \(\sigma\) 的和在 \(\upsilon=\ell\) 时等于4，其他三个 \(\upsilon\) 时均为0。下端要求字符和为 \(-64\)，因此 \(\operatorname{Tr}(f\ell)=1\)。引理5.3再以 \(\delta=0\) 迫 \(j=k=0\)。于是全部等号细化只能是



\[
Q_{f,g,\ell,m}=ad+bc+\operatorname{Tr}(ac)+fa+ga^2+\ell c+mc^2,
\qquad \operatorname{Tr}(f\ell)=1.\tag{\text{原稿 5.9}}
\]



反方向也要检查：这些映射在零频处 \(a_z=c_z=0\)，故三个原始相位均为0；上式的字符和确为 \(-64\)。所有极化保持非退化，原稿的计数恒等式于是给出



\[
v_4=M_4+\frac{2^4}{4}(-64)=4627-256=4371.
\]



因此列出的不只是必要候选，确实个个满足要求。计数时 \(f\) 有3个非零选择；每个 \(f\) 下有2个 \(\ell\) 使 \(f\ell\in\{\alpha,\alpha^2\}\)；\(g,m\) 各4种。因此共 \(3\cdot2\cdot4\cdot4=96\)。



\textbf{96是固定极化、固定坐标中的映射数，不是群同构类型数。} 全部 \(2^{16}\) 候选到96的缩减靠充要条件完成，不靠“程序没找到别的”。下一步才比较这96个之间的严格等价。



\Needspace{7\baselineskip}
\section{四次显式输入换基如何给出传递性}


\Needspace{5\baselineskip}
命题5.5要证明它们全都等价于



\[
Q_*=ad+bc+\operatorname{Tr}(ac)+a+\alpha c.
\]



原稿使用 \(\operatorname{SL}_2(R)\) 作用于列向量 \((x,y)^{\mathsf T}\)，\(x=a+\varepsilon b\)、\(y=c+\varepsilon d\)。对任意二阶矩阵，交替表达式 \(xy'+x'y\) 在换基后乘以行列式；行列式为1便同时保持式(5.3)的两个标量极化。下面每一步都是对完整映射作拉回，即计算 \(Q(M(x,y))\)，不是任意添加一个线性映射。



\Needspace{5\baselineskip}
第一步取



\[
M_1=\begin{pmatrix}1&\varepsilon m\\ \varepsilon g&1\end{pmatrix}.
\]



其行列式为 \(1+\varepsilon^2mg=1\)，所以可逆。它使 \(a,c\) 不变，而 \(b\mapsto b+mc\)、\(d\mapsto d+ga\)。基准项变成



\[
a(d+ga)+(b+mc)c+\operatorname{Tr}(ac)
=ad+bc+\operatorname{Tr}(ac)+ga^2+mc^2.
\]



这两个新项恰与原有 \(ga^2+mc^2\) 相消。于是 \(g,m\) 被消去。注意相消来自实际拉回的基准项；如果未经说明就“忽略线性项”，无法得到严格等价。



第二步取 \(M_2=\operatorname{diag}(f^{-1},f)\)。它将 \(a,b\) 乘 \(f^{-1}\)，\(c,d\) 乘 \(f\)。三个基准项均不变，\(fa\) 变成 \(a\)，\(\ell c\) 变成 \(f\ell c\)。迹一条件只剩 \(f\ell=\alpha\) 或 \(\alpha^2\)。



若已经为 \(\alpha\)，便得到 \(Q_*\)。若为 \(\alpha^2\)，第三步取



\[
M_3=\begin{pmatrix}1&\alpha\\0&1\end{pmatrix}.
\]



这给 \(a\mapsto a+\alpha c\)、\(b\mapsto b+\alpha d\)，另两坐标不变。\(ad+bc\) 的两个额外 \(\alpha cd\) 相消，但



\[
\operatorname{Tr}((a+\alpha c)c)
=\operatorname{Tr}(ac)+\alpha c^2+\alpha^2c.
\]



再加 \(a+\alpha c\) 及原来的 \(\alpha^2c\)，普通 \(c\) 项的系数成为 \(\alpha^2+\alpha+\alpha^2=\alpha\)，另留下 \(\alpha c^2\)。最后用



\[
M_4=\begin{pmatrix}1&\varepsilon\alpha\\0&1\end{pmatrix}
\]



使 \(b\mapsto b+\alpha c\)，基准项多出 \(\alpha c^2\)，恰好消掉残项。这四步的目标变换始终是 \(T=\operatorname{id}_W\)。



因此，对96个映射中的每一个，已经写出一条通向 \(Q_*\) 的可核查路径。这是传递性证明，不需要先枚举完整变换群，也不需要以“不变量相同”代替同构证明。



\Needspace{7\baselineskip}
\section{变换群、具体群与本章的逻辑终点}


原稿注5.6还说明固定两个标量极化的全体等距群就是 \(\operatorname{SL}_2(R)\)。保持两个极化的变换必与 \(A=B_1^{-1}B_2\) 交换，故是 \(R\) 线性的。若其行列式为 \(d\)，第一极化的保持要求 \(\lambda((d-1)r)=0\) 对所有 \(r\in R\) 成立；非退化配对迫 \(d=1\)。一个可作为可逆矩阵首列的列向量必须至少有一个单位坐标，故有 \(16^2-4^2=240\) 种；对每个首列，行列式映射满射到 \(R\)，每个纤维16元，于是群阶为 \(240\cdot16=3840\)。



这个群阶计算是补充核对，不承担前面的轨道穷尽性。即使再允许目标换基，也不能把已经同属一个较小群轨道的96个映射拆开。



最后，原稿式(5.11)给出不依赖抽象存在性的群乘法。取集合 \(K^4\times K\)，令



\[
F(v,v')=ad'+bc'+\operatorname{Tr}(ac')+(aa')^2+\alpha(cc')^2.
\]



它在二元意义下双线性；例如 \((aa')^2\) 对 \(a,a'\) 各自加法可加。取对角得 \(a^4=a\)、\(c^4=c\)，故 \(F(v,v)=Q_*(v)\)。按原稿式(2.1)定义 \((v,w)(v',w')=(v+v',w+w'+F(v,v'))\)，得到阶1024的 \(G_*\)。三个标量极化非退化，是因为 \(1,t,1+t\) 均为 \(R\) 的单位。因此命题2.1保证它确属要求的群类。



至此八维分类的四项任务都已闭合：所有下端极化纳入模型，全部细化纳入候选，充要条件筛出96个，显式严格换基将它们合并为一个群同构类型。这里的局部环、相位公式与Arf判据是已有工具；本文保留的研究对象是这些工具合成后的具体极值分类，而不是新的通用二次型理论。



\chapter{高维等号分类：相位常数、整个余块与共同中心}


本章对应原稿§6：引理6.1（基线第12页）、引理6.2、6.3（第13—14页）、定理1.2的完整证明及注6.4（第14页）。分类范围始终是偶数半维数 \(n=4+2r\ge4\)，其中 \(\dim_{\mathbb F_2}V=8+4r\)。\(r=0\) 是上一章的八维基础情形。



已经知道下端等号的极化有一个 \(p^2\) 块对及 \(r\) 个简单 \(p\) 块对。因此引理5.1给出共同正交分解 \(V=V_0\oplus V_1\)，\(\dim V_0=8\)、\(\dim V_1=4r\)，且完整平方映射自动分成 \(Q=Q^{(0)}\perp Q^{(1)}\)。本章真正需要证明的不是“极化可分解”，而是完整细化的限制：核心属于第五章的96个映射，整个余块是正型。不能从整体Gauss和为正直接得出这两条。



\Needspace{7\baselineskip}
\section{为什么每个块必须保留自己的原点相位}


\Needspace{5\baselineskip}
原稿式(5.2)把半单余块的全部细化写成



\[
Q^{(1)}=\sum_{\nu=1}^r
\left(u_\nu v_\nu+F_\nu u_\nu+G_\nu u_\nu^2
+L_\nu v_\nu+M_\nu v_\nu^2\right),
\]



所有系数在同一个 \(K=\mathbb F_4\) 中。它的极化只来自 \(u_\nu v_\nu\)。余下各项都是二元线性项，但其中平方项与普通一次项对 \(K\) 上齐次性的影响不同，暂时一项也不能删。



对标量分量 \(\operatorname{Tr}(zQ^{(1)})\)，取频率 \(\operatorname{Tr}\sum_\nu(\rho_\nu u_\nu+\sigma_\nu v_\nu)\)。一个平面上的指数是



\[
\operatorname{Tr}\left(
zuv+(\rho+zF+z^2G^2)u+(\sigma+zL+z^2M^2)v\right).
\]



先对 \(u\) 求和，只留下



\[
v=(\rho+zF+z^2G^2)/z,
\]



\Needspace{5\baselineskip}
并产生因子4。余下指数得到原稿式(6.1)：



\[
\mathcal W_{1,z}=4^r(-1)^{E_{1,z}},
\]



\[
E_{1,z}=\sum_\nu\operatorname{Tr}
\frac{(\rho_\nu+zF_\nu+z^2G_\nu^2)
(\sigma_\nu+zL_\nu+z^2M_\nu^2)}{z}.
\]



这也是CMP2018式(4)，p.313的直接使用；它允许非零的额外函数 \(g\)，所以原点常数不能被忽略。展开一个分子再除以 \(z\)，各类项为



\[
\begin{aligned}
&\rho\sigma/z,\quad \rho L,\quad z\rho M^2,
\quad F\sigma,\quad zFL,\quad z^2FM^2,\\
&zG^2\sigma,\quad z^2G^2L,\quad z^3G^2M^2.
\end{aligned}
\]



对三个 \(z\ne0\) 相加时，含 \(z^{-1},z,z^2\) 的和为零，常数项保留，\(z^3\) 项保留一次。再用迹对平方不变，便得



\[
D_1:=\sum_{z\ne0}E_{1,z}
=\operatorname{Tr}\sum_\nu
(L_\nu\rho_\nu+F_\nu\sigma_\nu+G_\nu M_\nu).
\tag{\text{原稿 6.2}}
\]



重要的第一项结论是：\(D_1\) 的极化恒为零。它可能是非零线性函数，也可能是常数1，并不自动恒零。



若后来证明了 \(F_\nu=L_\nu=0\)，则零频处每个 \(E_{1,z}(0)\) 都等于



\[
\delta=\operatorname{Tr}\sum_\nu G_\nu M_\nu,
\]



且 \(D_1\equiv\delta\)。因为共有三个相位，在 \(\mathbb F_2\) 中 \(\delta+\delta+\delta=\delta\)。若 \(\delta=1\)，整个余块的三个原始Gauss和都为 \(-4^r\)，并非“二元bent函数regular”就能使它们为正。regular描述Walsh表示的统一振幅，零频符号仍可负。



\Needspace{7\baselineskip}
\section{全局最大自逆数起初只给两个符号相等}


\Needspace{5\baselineskip}
对共同正交和，三个标量Walsh变换分别按块相乘：



\[
\mathcal W_z(\xi_0,\xi_1)
=16\,4^r(-1)^{E_{0,z}(\xi_0)+E_{1,z}(\xi_1)}.
\]



\Needspace{5\baselineskip}
整体属于最大自逆数分支，等价于三个整体原点相位均为零。因此只能先写



\[
E_{0,z}(0)+E_{1,z}(0)=0,
\qquad E_{0,z}(0)=E_{1,z}(0).\tag{\text{原稿 6.4}}
\]



此时每个 \(z\) 上都可能是 \(1+1=0\)；而且起初不能排除不同 \(z\) 的相位不同。这正是引理6.2要处理的补偿问题。



\Needspace{5\baselineskip}
定义未单独调整常数的三相位和



\[
D_0=\sum_zE_{0,z},\qquad D_1=\sum_zE_{1,z},\qquad D=D_0+D_1.
\]



由于整体原点相位为零，\(D\) 正是第三章所用的规范三对偶和，满足 \(D(0)=0\)。但两个限制块的 \(D_0(0)\)、\(D_1(0)\) 可以同时为1。若先把两者分别改成零，就改变了限制块实际Gauss符号，无法正确证明补偿被排除。



\Needspace{7\baselineskip}
\section{引理6.2的完整条件链}


\Needspace{5\baselineskip}
下端等号除最大自逆数外，还给出原稿式(6.3)：



\[
D|_{\ker C}=0,
\qquad \sum_{\xi\in V^*}(-1)^{D(\xi)}=-2^{2n-2}.
\]



\Needspace{5\baselineskip}
下面按必要顺序使用这两个条件。



\textbf{先看整个余块的根方向。} \(D_1\) 的极化为零，所以 \(V_1^*\) 全部在整体 \(C\) 的根中。对每个 \(\xi_1\)，根消失条件给



\[
0=D(0,\xi_1)=D_0(0)+D_1(\xi_1).
\]



因此 \(D_1\) 必须是常数，而不仅是字符和为零。式(6.2)与迹配对非退化迫每个 \(F_\nu=L_\nu=0\)。于是 \(D_1\equiv\delta\)，余块三个原始相位都为 \(\delta\)。由整体正号可知核心的三个原始相位也都为 \(\delta\)，故上一章式(5.7)中的 \(\gamma=\sum_zE_{0,z}(0)=\delta\)。



\textbf{再看核心的剩余根方向。} 式(5.6)中二次部分是 \(\operatorname{Tr}(\sigma\upsilon)\)，根为 \(\rho,\tau\) 方向。限制整体 \(D\) 到这些方向，得到



\[
D(\rho,0,\tau,0;0)
=\operatorname{Tr}(\theta\rho+\eta\tau)+\gamma+\delta
=\operatorname{Tr}(\theta\rho+\eta\tau).
\]



现在 \(\gamma+\delta=0\) 是已经证明的相等常数互相抵消，而不是人为在两块上各删一次常数。根消失于是迫 \(\theta=\eta=0\)。



\Needspace{5\baselineskip}
\textbf{然后才用下端的负号。} 此时整体函数简化为



\[
D=\operatorname{Tr}(\sigma\upsilon+f\upsilon+\ell\sigma).
\]



它与全部 \(V_1^*\) 坐标及核心 \(\rho,\tau\) 无关，故其字符和为



\[
2^{4r}\cdot64(-1)^{\operatorname{Tr}(f\ell)}.
\]



由于 \(n=4+2r\)，所需下端值为 \(-2^{2n-2}=-2^{6+4r}=-64\cdot2^{4r}\)。因此



\[
\operatorname{Tr}(f\ell)=1.
\]



这是整条证明中不可省略的负号输入。若只知道整体最大自逆数而不知道下端，\(\operatorname{Tr}(f\ell)\) 可以为0，下一步的各向异性便不能成立。



\textbf{最后回到原始核心符号。} 此时已知核心三个原始相位均为同一 \(\delta\)，同时 \(\eta=\theta=0\)、\(\operatorname{Tr}(f\ell)=1\)。上一章引理5.3适用，得到 \(j=k=0\)、\(\delta=0\)。所以结论不是假定而是推出：核心三个原始和均为正16，整个余块三个原始和均为正 \(4^r\)。其完整参数限制为



\[
\eta=\theta=j=k=0,\quad\operatorname{Tr}(f\ell)=1,
\quad F_\nu=L_\nu=0,
\quad\operatorname{Tr}\sum_\nu G_\nu M_\nu=0.
\]



\(r=0\) 时余块是零维空间，空和 \(\delta=0\)；上述公式仍成立，但不应虚构一个正维的余块或扩域。这也与第五章的独立八维证明相符。



引理排除的是“八维核心和整个半单余块都具有负的三个原始Gauss和，而整体因此为正”。它没有说在事先任意选好的余块平面分解中，每个小平面的Gauss和都正。后者是错误的更强说法。



\Needspace{7\baselineskip}
\section{正余块何时才真正成为四元域上的齐次二次型}


\Needspace{5\baselineskip}
上述筛选后，余块变为



\[
q(u,v)=\sum_\nu(u_\nu v_\nu+G_\nu u_\nu^2+M_\nu v_\nu^2),
\qquad\operatorname{Tr}\sum_\nu G_\nu M_\nu=0.
\tag{\text{原稿 6.5}}
\]



这时确实有 \(q(cu,cv)=c^2q(u,v)\) 对所有 \(c\in K\) 成立，极化



\[
B_K((u,v),(u',v'))=\sum_\nu(u_\nu v'_\nu+u'_\nu v_\nu)
\]



是 \(K\) 双线性。普通一次项已经由等号条件消去；并非从一开始便因出现 \(K\) 坐标而具有这个齐次性。核心中的 \(\operatorname{Tr}(ac)+fa+\ell c\) 一般不满足 \(K\) 齐次性，所以不能对核心直接套同一种域上二次型分类。



AFH p.509的有限域二次型分类以及Theorem 2.2、pp.511—512可以解释此处的型别；其基域取 \(K\)，正维定义域可通过一个 \(K\) 线性同构视为 \(\mathbb F_{4^{2r}}\)。但原稿还给出一个可自行重建的双曲平面归纳，避免把“由标准理论”当作黑箱。



\Needspace{7\baselineskip}
\section{正型余块的双曲分解归纳}


先设 \(r\ge1\)。式(6.1)在零频处给三个标量和都为 \(+4^r\)。利用零点指示函数，\(q=0\) 的向量数为



\[
z_q=\frac{4^{2r}+3\cdot4^r}{4}>1.
\]



因为零向量只占一个，所以存在非零 \(u\) 使 \(q(u)=0\)，称为非零迷向向量。非退化性保证存在 \(v\) 使 \(B_K(u,v)=1\)。若 \(v\) 还不是迷向，可以修正为



\[
v'=v+q(v)u.
\]



\Needspace{5\baselineskip}
齐次性与极化公式给



\[
\begin{aligned}
q(v')&=q(v)+q(v)^2q(u)+q(v)B_K(v,u)\\
&=q(v)+0+q(v)=0,
\end{aligned}
\]



且 \(B_K(u,v')=1\)。因此 \(u,v'\) 线性独立，其张成的 \(K\) 平面非退化，并且



\[
q(au+bv')=a^2q(u)+b^2q(v')+abB_K(u,v')=ab.
\]



这就是一个 \(H_+\) 平面。由于该平面非退化，全空间是它与 \(K\) 正交补的直和。正交补维数减少2个 \(K\) 维，极化仍非退化；每个标量Gauss和按正交和相乘，而已分出的平面三个因子都为 \(+4\)，所以正交补三个和都为 \(+4^{r-1}\)。当补空间非零时，再用相同零点计数找到新的非零迷向向量，重复上述构造。最终



\[
q\cong_K H_+^{\perp r},\qquad H_+(u,v)=uv.
\]



全过程只更换输入的 \(K\) 线性辛基，目标 \(K\) 一直没动。归纳步骤不是“整个正型所以每个原平面正型”：我们可以混合原坐标来选择新的迷向平面。这里的正型判定与双曲分解是经典理论，不是新的分类方法。



\Needspace{7\baselineskip}
\section{共同目标为什么使两块换基能够真正拼接}


第五章把核心96个映射化到 \(Q_*\) 的变换都满足 \(T=\operatorname{id}_W\)；上一节对余块的归一也满足 \(T=\operatorname{id}_W\)。因此若它们的输入变换分别为 \(S_0,S_1\)，块对角变换 \(S_0\oplus S_1\) 在整体上是合法的同一次严格等价。



若两块各自只知道存在某个目标变换 \(T_0,T_1\)，却没有保证二者相同，不能直接拼成整体同构。原因是 \(W\) 是同一个二维中心，不是“每块各有一个互不相关的中心”。本稿通过更强的目标固定归一，已经解决这个兼容问题，而没有把它留给不明确的“适当换基”。



原稿注6.4还检查不同中心粘合方式。\(\operatorname{GL}_2(\mathbb F_2)\) 的6个元素，在 \(K\) 记号中恰是



\[
T(w)=c w^\sigma,
\quad c\in K^\times,
\quad \sigma=\operatorname{id}\ \text{或}\ (w\mapsto w^2).
\]



这6个映射互异且均二元线性可逆，数目正好为 \((4-1)(4-2)=6\)。在一个正平面上，目标变化能由输入变化提升：



\[
H_+(cu^\sigma,v^\sigma)=c(uv)^\sigma=T(H_+(u,v)).
\]



因此给不同正因子粘合中心时出现的表面差别，可在该正因子输入中吸收，不要求核心再采用一个不兼容的目标变换。这是对中心识别的额外核查，并不意味着任意群类的所有中心积都只有一个粘合法。



\Needspace{7\baselineskip}
\section{两个负平面为什么仍然给正余块}


\Needspace{5\baselineskip}
考虑原稿注6.4的



\[
A_-(u,v)=uv+u^2+\alpha v^2.
\]



一个平面的 \(G=1,M=\alpha\)，故 \(\delta=\operatorname{Tr}(\alpha)=1\)，三个原始Gauss和均为 \(-4\)。但两个这样的平面的总 \(\delta=1+1=0\)，整个四维 \(K\) 余块是正型。



为了看到“混合”不是口号，原稿提供以下四列组成的 \(K\) 线性变换，旧坐标顺序为 \((u_1,v_1,u_2,v_2)\)：



\[
e=(0,1,0,1),\quad f=(0,1,1,1),\quad
g=(1,0,1,0),\quad h=(\alpha,1,\alpha,0).
\]



对 \(q=A_-\perp A_-\)，逐项代入得 \(q(e)=q(f)=q(g)=q(h)=0\)。例如



\[
q(h)=(\alpha+\alpha^2+\alpha)+\alpha^2=0.
\]



其极化仍为 \(u_1v'_1+u'_1v_1+u_2v'_2+u'_2v_2\)。直接得到 \(B_K(e,f)=B_K(g,h)=1\)，其余不同组的四个交叉配对均为0；例如 \(B_K(e,g)=1+1=0\)、\(B_K(f,h)=\alpha+\alpha=0\)。这说明四列线性独立，并组成双曲辛基。于是



\[
q(xe+yf+zg+wh)=xy+zw.
\]



所以这是 \(A_-\perp A_-\) 到 \(H_+\perp H_+\) 的显式目标固定等距。它并没有把每个原来的负平面单独变成正平面，而是混合了两个平面的坐标。由此 \(Q_*\perp A_-\perp A_-\) 仍然达到对应下界，且不产生新群类型。



\Needspace{7\baselineskip}
\section{定理1.2的必要性、充分性与群阶}


\Needspace{5\baselineskip}
现在按原稿证明逐一闭合三个条件的等价。



从数值条件出发：最大 \(s\) 与下端等号，经命题3.5迫出恰一对 \(p^2\) 核心及简单 \(p\) 余块，同时迫 \(D\) 根上消失、字符和为 \(-2^{2n-2}\)。引理5.1提供共同极化模型，引理6.2筛出核心96个细化及正余块。命题5.5与引理6.3以同一目标固定将二者归一，所以



\[
Q\sim Q_*\perp H_+^{\perp r}.
\]



反过来取这个标准形。每个非零标量极化都是非退化块的直和，所以仍非退化；不能只验证三个原始Gauss和而忘记群类条件。三个原始Gauss和为 \(16\cdot4^r=2^{4+2r}=2^n\)，故 \(s=2^{2n}+3\cdot2^n\)。其三对偶和的字符和按块相乘为



\[
(-64)\cdot16^r=-2^{6+4r}=-2^{2n-2}.
\]



于是第三Walsh矩 \(\mathcal P=2^{3n}\mathcal G(D)=-2^{5n-2}\)，再由式(3.2)得



\[
v_4=M_n+\frac{\mathcal P}{4\cdot2^{2n}}
=M_n-2^{3n-4}.
\]



最后解释群结构。\(H_+\) 对应 \(H(K)\)，在 \(K^2\times K\) 上的乘法为



\[
((u,v),w)((u',v'),w')
=((u+u',v+v'),w+w'+uv').
\]



其阶为 \(4^3=64\)，平方为 \(uv\)，中心为 \(K\)。由引理2.2，正交和对应沿同一个 \(W\) 的中心积，而不是普通直积。含核心和 \(r\) 个正因子时，外直积阶 \(1024\cdot64^r\)，需要商去中心中阶为 \(4^r\) 的核，故



\[
|G|=1024\cdot64^r/4^r
=2^{10+4r}=2^{2n+2}.
\]



命题2.1再把完整平方映射的严格等价转换为群同构，得到每个固定偶数 \(n\ge4\) 恰一个整个群同构类型。这里证明了标准分解存在及全群类型唯一，没有证明某个给定群中的八维核心子空间唯一，更没有分类所有中心分解。严格等价保留全部线性项，避免把只在交换子层面同族的群错误合并。



\chapter{两个范围例子：结论为什么不能再加强}


本章对应原稿§7，第15—16页：例7.1、引理7.2和例7.3。前一个例子说明“整体最大自逆元素数”不能单独推出“每个限制块都正”；后一个例子说明偶数半维数的八维核心分类不能直接推广到奇数半维数。这两个反例针对不同的更强命题，都不否定定理1.1或1.2。



\Needspace{7\baselineskip}
\section{十二维例子：两个负块可使整体全正}


令 \(K=\mathbb F_4\)、\(\alpha^2+\alpha+1=0\)，所有块均取值于同一个 \(K\)。原稿取八维二元核心



\[
Q_0^-(a,b,c,d)=ad+bc+\operatorname{Tr}(ac)+\alpha a^2+d^2,
\]



\Needspace{5\baselineskip}
以及四维二元平面



\[
A_-(u,v)=uv+u^2+\alpha v^2.
\]



两者的线性细化不改变极化；核心极化与第五章的 \(p^2\) 模型相同，平面极化与 \(H_+\) 相同。因此各自三个非零标量极化都非退化，共同正交和也满足原群类假设。



对核心而言，式(5.1)中的参数是 \(g=\alpha,j=1\)，其余全部为0。在引理5.3的记号中 \(e=jk=0\)、\(u=fj^2+\ell k^2=0\)、\(v=gj+mk=\alpha\)，故三个原点相位均为



\[
E_{0,z}(0)=\operatorname{Tr}(\alpha)=1.
\]



所以核心三个原始Gauss和均为 \(-16\)。平方映射的定义域有256个向量，其相应群自逆元素数为



\[
s(Q_0^-)=256+3(-16)=208.
\]



对平面，\(G=1,M=\alpha\)，\(\delta=\operatorname{Tr}(\alpha)=1\)，三个原始Gauss和均为 \(-4\)，所以



\[
s(A_-)=16+3(-4)=4.
\]



但共同正交和在每个标量分量上将两和相乘，得到 \((-16)(-4)=+64\)。此时 \(V\) 的二元维数为 \(8+4=12\)，半维数 \(n=6\)，群阶为 \(4\cdot2^{12}=16384\)，整体自逆元素数达到



\[
s=2^{12}+3\cdot64=4096+192=4288.
\]



也就是说，两个限制块都不在各自最大 \(s\) 分支上，但整体却在最大 \(s\) 分支上。零点数不会简单相乘，因为 \(Q_0^-(x)+A_-(y)=0\) 允许两个非零值互相抵消；真正按块相乘的是每个标量字符和。



这还没有检查是否达到四元子群下界。保留实际对偶的原点常数，由式(5.6)、式(6.2)可得



\[
D_0=\operatorname{Tr}(\sigma\upsilon)+1,
\qquad D_1=1.
\]



核心的实际三对偶字符和为 \(-64\)，平面为 \(-16\)；相加相位时两个常数1抵消，整体



\[
D=\operatorname{Tr}(\sigma\upsilon),
\qquad\mathcal G(D)=(-64)(-16)=+1024.
\]



下端本来要求 \(-2^{2n-2}=-1024\)，恰好相反。因此引理6.2的负号假设没有满足。本例不能被用作该引理或主分类的反例。



数值可逐项对照：\(h=64\)，



\[
M_6=\frac{64^4+9\cdot64^3+36\cdot64^2-36\cdot64+8}{24}
=803403.
\]



由第三章 \(v_4=M_n+(h/4)\mathcal G(D)\)，本例得到



\[
v_4=803403+16\cdot1024=819787.
\]



而真正的标准下端 \(Q_*\perp H_+\) 给 \(\mathcal G(D)=-64\cdot16=-1024\)，所以



\[
v_4=803403-16384=787019.
\]



这两个映射的极化类型相同，整体自逆元素数相同，却因完整线性细化不同而给不同的四元子群数。它再次说明交换子结构不足以替代完整平方映射。本例也清楚标出补偿排除证明真正使用下端条件的地方：负号迫 \(\operatorname{Tr}(f\ell)=1\)，而本例 \(f=\ell=0\)，没有各向异性障碍。



另一方面，\(Q_*\perp A_-\perp A_-\) 确能达到更高维相应下端。这并不矛盾：整个余块 \(A_-\perp A_-\) 的三个原始和为正16，符合引理6.2；只是原来选的两个小平面为负。第六章已经给出它们混合成两个正平面的显式换基。这里不能把“核心与整个余块”偷换成“所有初始小平面”。



\Needspace{7\baselineskip}
\section{目标换基怎样改变极化算子}


为了证明下一个例子真的不存在八维核心，而非“没有在当前坐标中看到”，必须允许全部合法的目标换基。令原有序标量极化为 \(B_1,B_2\)。目标换基后可写成



\[
B'_1=aB_1+bB_2=B_1(aI+bA),
\qquad B'_2=cB_1+dB_2=B_1(cI+dA),
\]



其中 \(\begin{pmatrix}a&b\\c&d\end{pmatrix}\in\operatorname{GL}_2(\mathbb F_2)\)。非零的第一行对应原来的某个非零标量极化，所以 \(aI+bA\) 可逆。于是



\[
A'=(B'_1)^{-1}B'_2
=(aI+bA)^{-1}(cI+dA).\tag{\text{原稿 7.1}}
\]



\Needspace{5\baselineskip}
另作输入换基只会将该算子相似变换。这是一种分式线性变换，而不是只允许交换两个块。



算子“半单”在这里是指扩到代数闭包后可对角化，等价于最小多项式没有重复不可约因子。若 \(A\) 半单，在一个对角化基下 \(A'\) 也是将每个特征值代入上述分式所得的对角矩阵，因此仍半单。即使不同特征值代入后相同，矩阵仍然对角化。反方向用逆目标换基得到，所以半单性是双向保持的，而不只是某些生成操作下保持。



这就是原稿引理7.2的用途：以完整允许变换下的不变量，排除隐藏的非半单核心。仅证明固定有序极化下的障碍，不足以处理群同构允许的目标变化；此引理补上了这一步。



\Needspace{7\baselineskip}
\section{十维奇数分支例子：确实达下端，却没有八维核心}


取 \(L=\mathbb F_8=\mathbb F_2[t]/(t^3+t+1)\)，\(d=t^2+t+1\)，绝对迹记为 \(\operatorname{Tr}_8\)。在 \(V=L^2\oplus K^2\) 上取同一个二元目标 \(W=\mathbb F_2^2\)，定义



\[
\begin{aligned}
Q_{\mathrm{odd}}(x,y,u,v)=\bigl(&\operatorname{Tr}_8(xy)+\operatorname{Tr}(uv),\\
&\operatorname{Tr}_8(txy+x+dy)+\operatorname{Tr}(\alpha uv)\bigr).
\end{aligned}
\]



它是第四章引理4.2的负端六维块与引理4.1的正四维块的共同目标正交和。六维块的三个标量极化系数为 \(1,t,1+t\)，均非零；四维块的三个系数也非零。因此整体满足要求的所有标量非退化假设，并对应真实的semi-extraspecial群。



已有块计算给六维原Gauss和 \((8,8,8)\)、三对偶字符和 \(-8\)；四维块给 \((4,4,4)\)、三对偶字符和16。因此整体三个原始和为 \(32\)、三对偶字符和为 \(-128\)。此时 \(\dim_{\mathbb F_2}V=10\)，\(n=5\)，\(h=32\)，\(N=1024\)，故



\[
|G|=4096,\qquad s=1024+96=1120.
\]



又 \(M_5=57467\)，奇数分支下端修正量是 \(2^{3n-5}=2^{10}=1024\)。直接换算为



\[
v_4=M_5+\frac{32}{4}(-128)=57467-1024=56443.
\]



\Needspace{5\baselineskip}
因此它确实满足奇数分支的最大自逆元素数和下界等号；不是由于放宽了数值条件才绕过八维结论。



为何没有 \(Q_*\) 因子？在六维部分，以迹配对识别算子，\(B_1\) 对应乘1、\(B_2\) 对应乘 \(t\)，所以 \(A\) 在两个 \(L\) 坐标上都为乘 \(t\)，其两个循环块的最小多项式是简单不可约三次多项式 \(t^3+t+1\)。在四维部分，\(A\) 为乘 \(\alpha\)，给简单 \(p\) 块对。因此整体 \(A\) 半单。



八维核心 \(Q_*\) 的 \(A\) 则有 \(p^2\) 块对，故非半单。若存在某个共同正交因子严格等价于 \(Q_*\)，即使先作任意允许的输入、目标换基，整体算子仍半单。共同正交分解使两个极化矩阵同时成为块对角，其比值 \(A\) 也块对角，核心上的限制就是整体 \(A\) 在一个不变子空间上的限制。但半单算子的限制仍半单：限制的最小多项式整除整体那个无重复因子的最小多项式。这与核心必须有 \(p^2\) 最小因子矛盾。



所以这个障碍排除全部允许目标换基后的核心，不只是当前写法中的核心。另一个简单维数提醒是：十维若硬减八维会余二维，而本题二维空间不能有三个非零标量极化都非退化；但上述半单性证明更强，它不要求剩余因子仍属于同一非退化群类，也已经排除了任意共同正交剩余部分。



这个例子说明统一的“所有下端都有 \(Q_*\)”猜想为假。它没有完成奇数分支完整平方映射的分类。命题3.5(a)已经限制奇数端点的极化类型，却还没有像偶数分支那样把所有线性细化及严格等价轨道分类完。本稿准确停在“奇数数值界及端点可达、偶数下端全分类”，不把前者升级为后者。



\chapter{已有工具、具体结论与两个保留附录}


本章对应原稿§8（基线第16—17页）和附录A、B（第17—18页）。这一部分有两种任务：说明本文结论与已有理论的实际接口；把正文使用的共同坐标和一个特定版本文献说明展开。它不是重新进行无边界文献审查，也不把“没有在已读文献找到完全相同表述”当作原创性证明。



\Needspace{7\baselineskip}
\section{哪些工具已经存在，本文究竟把什么合成起来}


\Needspace{5\baselineskip}
可以将已读文献的作用按证明位置理解。



\begin{longtable}{@{}p{0.3033\linewidth}@{\hspace{.025\linewidth}}p{0.3033\linewidth}@{\hspace{.025\linewidth}}p{0.3033\linewidth}@{}}
\toprule
来源及已核位置 & 在本稿中的实际作用 & 不能由它直接声称的结论 \\
\midrule
\endhead
Kaur–Kulshrestha，arXiv:1210.3790v1，Lemma 2.4，PDF第6—7页 & 正交和与沿共同中心的中心积对应 & 不是本稿四元子群极值分类 \\
Kaur–Kulshrestha，arXiv:1510.06583v1，Proposition 3.5、Remark 3.6，第8页 & 此群类的不可约字符及三个非线性字符的来源 & 没有以此自动核读Zahinda原文的全部分类 \\
Abdukhalikov–Feng–Ho（AFH），Theorem 2.1，印刷pp.509—510 & 实际二元标量Walsh对偶；允许二元线性细化 & 不是三个分量的共同严格等价分类 \\
AFH，p.509及Theorem 2.2，pp.511—512 & 已满足较大基域齐次性后的型别、相位及正余块解释 & 不能因使用四元域坐标便将核心当作四元域齐次型 \\
Çeşmelioğlu–Meidl–Pott（CMP），式(4)，p.313 & 一般Maiorana–McFarland对偶公式，含任意额外函数及原点相位 & 不是不经联立等号条件便自动得到96个筛选 \\
CMP，Theorem 2(i)，p.314 & 标准迹乘积正端构造，须保留子域线性及相应假设 & 不把任意二元线性置换误作四元域线性置换 \\
Wang–Fu–Wei（WFW），arXiv:2301.00581v1，Theorem 6，式(35)，pp.23—24 & \(D\equiv0\) 上端的部分差集参数 & 不直接覆盖非半单下端 \\
交替矩阵笔及局部等价理论，LM、PFMR、MPG的指定版本和位置 & 成对循环块、局部环输入变换的框架 & 同族不等于群同构；不能丢弃完整线性细化 \\
\bottomrule
\end{longtable}


AFH与CMP的全文已实际取得并完成项目中的命题级比较，不是当前取文缺口。它们不是“只与关键词有关”的背景文献：正文第五、六章的实际相位正是旧公式的代入；给定等号极化后，96个筛选及核心—整个余块的补偿排除，都是旧公式加经典各向异性判据的短推导。本稿不把这些步骤包装成新的通用对偶理论、规范形算法或新方法。



当前具体合成结论是：先由群计数得到三阶Walsh矩；利用交替笔的维数亏损求出最优数值界及全部等号条件；在偶数 \(n\ge4\) 下联立完整细化条件，再以同一目标固定归一，得到整个下端群的唯一同构类型。已经核读的上述定理未直接代替整条结论，这不意味着其他文献不可能已有同一特例。历史位置与数学证明的有效性是两个问题，不能混写。



Zahinda关于totally orthogonal特殊2群的完整细化分类是否已经覆盖中心 \(C_2^2\) 的有关结构，仍是具体历史问题；本文使用其Arf判据的引用链已经明确，未据此声称核读该原文所有结果。CMP2021、Fitzgerald–Yucas、Bornand及后续结构论文也仅保留真实未核的具体范围，不把每个未读书目自动改写为本稿数学证明缺口。细节留在项目文献记录，不需要塞进数学证明正文。



\Needspace{7\baselineskip}
\section{最大自逆数分支与totally orthogonal的准确联系}


这段在正文§8中给出另一种文献入口，并不承担最优界的自足证明。不熟悉表示论的读者可先理解结论：这里“最大自逆元素数”恰是这个群类的“所有不可约字符Frobenius–Schur指标均为正”分支。它使Zahinda文献与本问题有真实联系，而非只因标题相似。



简要补充必要记号。一个复表示把群元映成可逆复矩阵，其字符 \(\chi(g)\) 为矩阵迹。不可约表示是不能分成更小不变子空间直和的表示。Frobenius–Schur指标定义为



\[
\nu(\chi)=\frac1{|G|}\sum_{g\in G}\chi(g^2).
\]



经典判据说明：不可约字符的指标为 \(+1\) 时可由实正交表示实现；所有不可约字符都满足这一性质便称totally orthogonal。这个表示论判据是已有背景，不在本文重新证明。



在本群类中，\(G/G'=V\) 是阶 \(N\) 的二元初等阿贝尔群，所以有 \(N\) 个线性字符，它们只取 \(\pm1\)，且 \(\chi(g^2)=1\)，指标均为1。按KK15 Proposition 3.5、Remark 3.6，还有三个度为 \(h=2^n\) 的不可约字符 \(\chi_\lambda\)，分别来自三个extraspecial商。其限制到中心为



\[
\chi_\lambda(w)=h(-1)^{\lambda(w)},\qquad
\lambda\in W^*\setminus\{0\}.
\]



这已经列完所有不可约字符：它们的度数平方和为 \(N+3h^2=4N=|G|\)。因为每个 \(g^2\in W\)，且每个 \(v\in V\) 有四个提升，指标变为



\[
\begin{aligned}
\nu(\chi_\lambda)
&=\frac1{4N}\sum_g h(-1)^{\lambda(g^2)}\\
&=\frac{4h}{4N}\sum_{v\in V}(-1)^{\lambda Q(v)}
=\frac{\mathcal G(\lambda Q)}h.
\end{aligned}
\]



所以这三个指标全为1，当且仅当三个原始Gauss和全为 \(+h\)，也就当且仅当 \(s=N+3h\)。KK12 Theorem 4.1、第12页重述了相应Arf判据，并将其归于Zahinda Theorem 3.5；这里准确地通过KK使用这条归属，没有将间接引用升级为已检查Zahinda全部原文。



\Needspace{7\baselineskip}
\section{三个对偶的闭合：\texorpdfstring{\(D\equiv0\)}{Dequiv0}的真正含义}


设三个原分量为 \(q_1,q_2,q_3=q_1+q_2\)。实际Walsh对偶由



\[
\mathcal W_{q_i}(\xi)=h(-1)^{q_i^*(\xi)}
\]



唯一决定，包括原点相位。在最大 \(s\) 分支上三者的原点相位都为0。三个对偶均为非零且互异的函数：每个对偶仍bent，零函数不是正维bent；若两个对偶相同，两个Walsh变换相同，由Fourier反演原分量相同，这将使第三标量分量为零，与极化非退化矛盾。



CMP Definition 1、p.309中的vectorial dual-bent，要求非零分量的实际对偶加零函数构成相同输出维数的函数向量空间。本题输出维数是2，非零元素只有三个，故这个要求恰为



\[
\{0,q_1^*,q_2^*,q_3^*\}\text{对加法封闭}
\quad\Longleftrightarrow\quad q_1^*+q_2^*=q_3^*
\quad\Longleftrightarrow\quad D\equiv0.
\]



在一般输出维数中，向量对偶的标签选择要另外处理；此处三个非零标签的任意重排都来自 \(\operatorname{GL}_2(\mathbb F_2)\)，因为该群有6个元素并忠实作用为三个非零向量的所有置换。因此不能把一般维数的复杂标签问题，照搬成本题尚未解决的障碍。



\Needspace{5\baselineskip}
有一条很容易误用的矩阵恒等式：



\[
B_1^{-1}+B_2^{-1}
=B_1^{-1}(B_1+B_2)B_2^{-1}.
\]



右边三个因子均可逆，所以左边可逆。这意味着任意两个对偶之和仍是非退化二元二次函数，因而bent。但它没有说明这个和就是第三个对偶；后者还要求完整二次细化相同，恰是 \(D\equiv0\)。由“每两个对偶的和bent”跳到“对偶闭合”，遗漏了核心问题。



还必须区分三件事：\(D\equiv0\) 是完整函数恒零；\(C=0\) 是 \(D\) 的极化恒零，只能推出 \(D\) 为线性函数（在整体规范 \(D(0)=0\) 下）；\(\mathcal G(D)=0\) 是整数字符和相消，可能来自非零线性根限制。若 \(D\equiv0\)，反而 \(\mathcal G(D)=N\)，不是0。



\Needspace{7\baselineskip}
\section{为什么\texorpdfstring{\(D\equiv0\)}{Dequiv0}上端明确属于旧定理的直接推论}


这一步不是笼统地说“与部分差集有关”，而是原稿§8已经实际完成的参数代入。最大 \(s\) 加 \(D\equiv0\) 提供同标签向量对偶，满足WFW Theorem 6的Condition A。取其参数为 \(p=2\)、输入维数 \(2n\)、输出维数2、regular系数 \(\epsilon=1\)、\(F=Q\)、\(F(0)=0\)，并取值集 \(I=\{0\}\)。输出维数条件 \(2\le n\) 在本题有效维数内成立。



旧定理给 \(E=Q^{-1}(0)\setminus\{0\}\) 为部分差集。需要用到的参数为



\[
z=\frac{N+3h}{4},\quad k=|E|=z-1,
\quad\lambda=\frac N{16}+\frac{3h}{4}-2.
\]



部分差集参数 \(\lambda\) 的含义是：对每个固定 \(x\in E\)，有 \(\lambda\) 个有序对 \((y,z')\in E^2\) 使 \(y-z'=x\)。二元加法中减号与加号相同，因此这些恰是 \(x,y,x+y\) 全在非零零点集中的计数。总计 \(k\lambda\) 个有序基，每个全零二维子空间有6个有序基，所以 \(6L=k\lambda\)。再用引理2.3的提升公式，



\[
v_4=16L+6z-5=\frac83 k\lambda+6z-5.
\]



为了看清与 \(M_n\) 的关系，写 \(N=h^2\)，代入可得



\[
\begin{aligned}
v_4
&=\frac{(h^2+3h-4)(h^2+12h-32)}{24}
+\frac{36h^2+108h-120}{24}\\
&=\frac{h^4+15h^3+36h^2-36h+8}{24}
=M_n+\frac{h^3}{4}.
\end{aligned}
\]



即 \(M_n+2^{3n-2}\)。这是明确的旧定理直接推论，应在贡献叙述中承认。它不同时给出非 \(D\equiv0\) 的下端，也不能将“regular”误当每个正交限制块的零频和为正。



\Needspace{7\baselineskip}
\section{EA等价与completed class为什么不能代替群同构}


本稿严格等价要求 \(Q'(Sx)=TQ(x)\)，其中 \(S,T\) 同时作用于完整映射，且二元线性可逆。CMP的extended affine（EA）等价更宽，允许



\[
\widetilde Q(x)=TQ(Sx+a)+Lx+b.
\]



输入平移 \(a\)、任意线性添加 \(L\) 及目标常数 \(b\) 不属于本稿群同构关系。它们也可能改变零点计数和实际对偶闭合。比如 \(uv\) 与 \(uv+u\) 仅差一次项，却分别有69和53个四元子群；第五章的基准核心与 \(Q_*\) 也同属一个EA类，却给不同的 \(v_4\)。



不过，不能仅以“下端对象自身不闭合”便排除全部EA路径，因为EA确实会改变闭合性。正确的不变量是极化秩：在上述全部EA操作下，三个极化满足



\[
\widetilde B_\lambda=S^{\mathsf T}B_{T^*\lambda}S.
\]



平移、添加线性项和常数不改变二次极化；目标 \(T\) 只置换三个非零标签。因此三逆极化和变为



\[
\widetilde C=S^{-1}CS^{-\mathsf T}.
\]



所以 \(\operatorname{rank}C\) 在完整EA范围内保持。下端 \(\operatorname{rank}C=4\)，而二次对偶闭合源有 \(C=0\)，故下端不能整体EA等价于这样的闭合源。这一论证针对完整EA范围，而非只证明严格等价不变量后越界使用。



CMP Theorem 1、p.311的前提是存在向量对偶，其结论是completed class中的\textbf{标量函数集合}取对偶后的关系。它允许分别找到三个标量的实现，不会保证它们由同一个输入变换、同一个线性添加和同一个目标操作同时实现。各标量分别可化到标准二次型，是比完整向量映射共同分类弱得多的信息。上述秩论证排除了共同闭合源的EA代表，并没有否认每个标量分别属于某个completed class。



LM Theorem 1.2计数的是同族类，而非本稿完整平方映射同构类；其另外一些分解定理还带指数 \(p\) 假设，本题指数为4，不能未经核对套用。PFMR在特征2的局部理论则确实保留对角平方所表示的二元线性项，因此不能以“齐次二次型”四字排除其覆盖。稿件保留显式换基，使证明不依赖一条未经检查的宽泛算法声明。



\Needspace{7\baselineskip}
\section{附录A：为什么三个核心分量可用同一套\texorpdfstring{\(\mathbb F_{16}\)}{mathbb F16}坐标}


正文引理5.2已经直接求和；附录A另给CMP公式(4)的实际参数对应。这样引用不是“看到同一种方法”，而是可逐项核验的同一个函数。



\Needspace{5\baselineskip}
取



\[
E=K[\beta],\qquad\beta^2+\beta=\alpha.
\]



因为 \(\operatorname{Tr}_{K/\mathbb F_2}(\alpha)=1\)，方程 \(w^2+w=\alpha\) 在 \(K\) 无解，所以 \(E\) 是 \(K\) 的二次扩域，即 \(\mathbb F_{16}\)。它不同于第五章有幂零元的环 \(R\)。四元域上的非平凡共轭使 \(\beta\mapsto\beta+1\)，故相对迹 \(\operatorname{Tr}_{E/K}(\beta)=1\)，而 \(\operatorname{Tr}_{E/K}(c)=0\) 对所有 \(c\in K\)。若 \(r+s\beta\in E\)，其相对迹就是 \(s\)。



\Needspace{5\baselineskip}
作一次共同输入坐标及频率坐标变化：



\[
X=(b+d)+\beta d,\quad Y=a+\beta c,
\qquad U=\upsilon+\beta\sigma,
\quad V=(\tau+\rho)+\beta\rho.
\]



每个变化均可逆，例如从 \(X=x_0+\beta x_1\) 恢复 \(d=x_1,b=x_0+x_1\)；从 \(Y\) 恢复 \(a,c\)。它们对三个分量相同，没有给每个标量偷偷选不同坐标。



\Needspace{5\baselineskip}
乘法公式为



\[
(r+\beta s)(u+\beta v)
=(ru+\alpha sv)+\beta(rv+su+sv).
\]



因此 \(XY\) 的 \(\beta\) 系数为 \((b+d)c+da+dc=bc+ad\)，从而



\[
\operatorname{Tr}_{E/K}(XY)=ad+bc.
\]



同样 \(UX\) 的相对迹为 \(\upsilon d+\sigma b\)，\(VY\) 的相对迹为 \(\tau c+\rho a\)，故



\[
\operatorname{Tr}_{E/\mathbb F_2}(UX+VY)
=\operatorname{Tr}(\rho a+\sigma b+\tau c+\upsilon d).
\]



\Needspace{5\baselineskip}
这精确运输了原来的Fourier配对，而不是暗中更换内积。



\Needspace{5\baselineskip}
记



\[
\mathcal H(Y)=\operatorname{Tr}(ac)+fa+ga^2+\ell c+mc^2,
\quad Y=a+\beta c.
\]



它是一个 \(K\) 值函数，不必为 \(K\) 齐次二次型。对全部八系数细化，逐分量定义



\[
\begin{aligned}
h_z&=(z\theta+z^2j^2)+\beta(z\eta+z^2k^2),\\
\pi_z(Y)&=zY+h_z,
\qquad g_z(Y)=z\beta\mathcal H(Y).
\end{aligned}
\]



\(z\ne0\)，所以 \(\pi_z\) 是 \(E\) 的仿射置换。由上面的乘法，\(\operatorname{Tr}_{E/K}(Xh_z)\) 恰为 \(b(z\eta+z^2k^2)+d(z\theta+z^2j^2)\)。再用迹对平方的恒等式以及 \(\operatorname{Tr}_{E/K}(\beta\mathcal H)=\mathcal H\)，得到严格恒等式



\[
\operatorname{Tr}(zQ^{(0)})
=\operatorname{Tr}_{E/\mathbb F_2}
\left(X\pi_z(Y)+g_z(Y)\right).
\]



CMP式(4)、p.313在本次\textbf{标量}使用中的参数是 \(p=2,m=4,s=1\)、分量参数1。这里 \(z\) 已经用来标记原三个分量，不能再把它与文献中的另一个标量参数混用。原式允许任意 \(g\)；负号因特征2消失。先对 \(X\) 求和，即得实际对偶相位



\[
\operatorname{Tr}_{E/\mathbb F_2}
\left(V\pi_z^{-1}(U)+g_z(\pi_z^{-1}(U))\right).
\]



\Needspace{5\baselineskip}
现在



\[
\pi_z^{-1}(U)=a_z+\beta c_z,
\]



\[
a_z=\upsilon/z+\theta+zj^2,
\qquad c_z=\sigma/z+\eta+zk^2.
\]



第一项相对迹给 \(\rho a_z+\tau c_z\)；第二项给 \(z\mathcal H(a_z+\beta c_z)\)。再把平方项移到系数，正得到原稿式(5.5)，连同 \(E_z(0)\) 一起保留。这一代入适用于全部 \(2^{16}\) 个细化，不仅适用于筛选后的96个。



当 \(\eta=\theta=j=k=0\) 时，整个 \(K\) 值映射还可共同写成



\[
Q^{(0)}=\operatorname{Tr}_{E/K}
\left(XY+\beta\mathcal H(Y)\right).
\]



因此核心确实属于一般 \(g\ne0\) 的MM表示，而非与MM完全无关。但CMP Theorem 2的后续闭合结论另取 \(g=0\)，不能从这个共同表示跳过联合筛选。此处正确区分的是“旧公式确实给出所需相位”与“旧定理是否自动给完整极值分类”，不能取其中一端代替另一端。



\Needspace{7\baselineskip}
\section{附录B：MPG特定印刷陈述的最小实例}


本附录仅针对项目保存的MPG13 \textbf{19页ePrint版本，第14页Proposition 7}。该版本印刷陈述把vanishing pencil归约到特征 \((0,0,H^m,0)\)，\(m\ge d/2\)。这里不把其他版本同编号的自同构命题混为一谈，也不据此否定整篇算法。正文主定理的自足证明不使用这个归约。



为理解例子，只需一个更小局部环：\(R_0=\mathbb F_2[H]/(H^2)\)，局部长度 \(d=2\)，\(t=1+H\)。它不是前面 \(p(t)^2\) 的八维核心，也不属于本稿semi-extraspecial群类。令 \(x=a+Hb,y=c+Hd\)，\(a,b,c,d\in\mathbb F_2\)，并令 \(\lambda(a+Hb)=b\)。取



\[
q_\infty=ad+bc,\qquad
q_0=ad+bc+ac+a^2+c^2.
\]



注意二元函数中 \(a^2=a\)、\(c^2=c\)；这些对角平方项没有二次极化，却是完整函数的重要部分。把每个环坐标按 \((b,a)\)、\((d,c)\) 排列，乘 \(t=1+H\) 的矩阵与配对矩阵为



\[
M=\begin{pmatrix}1&1\\0&1\end{pmatrix},\qquad
T=\begin{pmatrix}0&1\\1&0\end{pmatrix},\qquad
TM=\begin{pmatrix}0&1\\1&1\end{pmatrix}.
\]



\(T,TM\) 都可逆，所以两个标量极化都非退化；\(M\) 的最小多项式和特征多项式同为 \((X+1)^2\)，所以是所需的循环正则primary笔。\(q_\infty\) 没有对角平方项，\(q_0\) 的两个对角特征为 \(H,H\)，完整特征为 \((0,0,H,H)\)。



\Needspace{5\baselineskip}
还需要核查原文vanishing条件，而不能只给一个任意非退化笔。可逆二元输入变化



\[
\phi(a,b,c,d)=(a,b+c,c,d+a+c)
\]



\Needspace{5\baselineskip}
会交换两个完整函数。直接代入：



\[
q_\infty(\phi)=a(d+a+c)+(b+c)c
=ad+bc+ac+a^2+c^2=q_0.
\]



由于 \(a,c\) 没变，额外的 \(ac+a^2+c^2\) 再加一次便相消，所以 \(q_0(\phi)=q_\infty\)；并且 \(\phi^2=\operatorname{id}\)，可逆性也已检查。因此两个有序笔都满足semi-vanishing，正是该印刷陈述采用的vanishing条件。



接下来不需要依赖某个未证明的迹不变量，只数完整函数的共同零点。源映射同时为零先要求两式相加为零：



\[
ac+a+c=0.
\]



将 \((a,c)=(0,0),(0,1),(1,0),(1,1)\) 代入，只有 \((0,0)\) 成立。此时 \(ad+bc=0\) 对 \(b,d\) 没有约束，故共同零点恰为4个。



印刷目标在 \(d=2,m\ge1\) 时，\(H^m\) 不是 \(H\) 就是0，因此目标第二式分别为



\[
ad+bc+ac+a^2,
\qquad ad+bc+ac,
\]



第一式仍为 \(ad+bc\)。在第一个目标，额外条件为 \(ac+a=0\)，允许 \((a,c)=(0,0),(0,1),(1,1)\)；这三种下，\(ad+bc=0\) 分别给4、2、2个 \((b,d)\)，合计8个共同零点。在第二个目标，额外条件为 \(ac=0\)，允许 \((0,0),(0,1),(1,0)\)，同样给4、2、2个，合计8个。



可逆输入换基是定义域的双射，保持共同零点数；再允许一个共同可逆目标换基，也保持“两个分量都为零”的条件。所以源的4个共同零点不可能变成目标的8个。该固定原域的同构陈述不包含扩域或改变函数的任意仿射操作，不能以那些更大的关系消除矛盾。



由此得到的是对\textbf{所保存版本和指定命题印刷陈述}的反例。它没有说其他版本相同编号的命题也错，没有检验整个算法，也不是本稿主定理原创的证据。19页版、17页会议存档和作者上传另一版的编号差别，已在项目第九轮专项记录中分别保存。公开论文保留此附录，是让特定文献说明具有可核查依据；读者可以完全不使用它来验证主定理1.1和1.2。



\Needspace{7\baselineskip}
\section{本章使用的可核对文献入口}


以下列出本章直接用到的正式书目信息及版本入口。链接用于核对来源，不表示这些第三方全文属于本项目发布包，也不授予重新分发权。



\begin{itemize}


\item K. Abdukhalikov, R. Feng, D. Ho, \emph{Quadratic bent functions and their duals}, Applicable Algebra in Engineering, \textbf{Communication} and Computing 35 (2024), 507—523. \href{https://doi.org/10.1007/s00200-022-00564-5}{DOI:10.1007/s00200-022-00564-5}。刊名Communication为单数。

\item A. Çeşmelioğlu, W. Meidl, A. Pott, \emph{Vectorial bent functions and their duals}, Linear Algebra and its Applications 548 (2018), 305—320. \href{https://doi.org/10.1016/j.laa.2018.03.016}{DOI:10.1016/j.laa.2018.03.016}。

\item D. Kaur, A. Kulshrestha, \emph{Strongly real special 2-groups}, \href{https://arxiv.org/abs/1210.3790v1}{arXiv:1210.3790v1}。中心积对应按原始v1 PDF为Lemma 2.4，不能以其他呈现方式的编号替换。

\item D. Kaur, A. Kulshrestha, \emph{Characters of real special 2-groups}, \href{https://arxiv.org/abs/1510.06583v1}{arXiv:1510.06583v1}。

\item M. L. Lewis, J. Maglione, \emph{Enumerating isoclinism classes of semi-extraspecial groups}, Proceedings of the Edinburgh Mathematical Society 63 (2020), 426—442. \href{https://doi.org/10.1017/S0013091519000531}{DOI:10.1017/S0013091519000531}；本稿定理编号按\href{https://arxiv.org/abs/1806.10511v2}{arXiv:1806.10511v2}。

\item J. Plût, P.-A. Fouque, G. Macario-Rat, \emph{Solving the “Isomorphism of Polynomials with Two Secrets” Problem for All Pairs of Quadratic Forms}, \href{https://arxiv.org/abs/1406.3163v3}{arXiv:1406.3163v3}，特别§4、PDF pp.15—18，Propositions 4.5、4.7。

\item G. Macario-Rat, J. Plût, H. Gilbert, \emph{New Insight into the Isomorphism of Polynomial Problem IP1S and Its Use in Cryptography}, \href{https://eprint.iacr.org/2013/799.pdf}{IACR ePrint 2013/799}。本稿附录B的对象是项目已保存19页版本，不将在线链接未来替换后的内容自动认作相同版本。

\item J. Wang, F.-W. Fu, Y. Wei, \emph{Bent Partitions, Vectorial Dual-Bent Functions and Partial Difference Sets}, \href{https://arxiv.org/abs/2301.00581v1}{arXiv:2301.00581v1}。

\end{itemize}


本解析是辅助自行阅读的解释材料。完成这些文件，不表示作者已经读完、核验全部证明或确认原创性；内部AI核查也不等同外部同行评审。若解析与对应正式论文实际定义或推导不符，应记录具体位置交由维护稿件的主对话处理，而不在另一答疑对话中无记录地改写论文。


\end{document}
